\documentclass[a4paper,11pt,twoside]{article}

% =========================================================================
%  Ecuaciones Diferenciales — Tema 5: EDO lineales de segundo orden con
%  coeficientes variables.  Parte I: puntos ordinarios y series enteras.
%
%  Compilación:   latexmk -pdf main.tex      (pdfLaTeX + biber)
%  Datos de las gráficas: python3 scripts/generar_datos.py  (ya incluidos)
% =========================================================================
% =========================================================================
%  Ecuaciones Diferenciales — Tema 5: EDO lineales de 2.º orden con
%  coeficientes variables. Parte I: puntos ordinarios y series enteras.
%
%  Estilo: artículo matemático sobre fondo oscuro (texto claro en Computer
%  Modern, como en las láminas de análisis complejo de referencia), con
%  enunciados numerados por sección, cajas con pestaña de título para
%  definiciones y teoremas y figuras TikZ/pgfplots adaptadas al fondo negro.
% =========================================================================

% --- IDIOMA Y CODIFICACIÓN ------------------------------------------------
% Se compila con LuaLaTeX (memoria dinámica para las superficies 3D);
% con pdfLaTeX funciona igualmente reduciendo la resolución de esas figuras.
\usepackage{iftex}
\ifPDFTeX
  \usepackage[utf8]{inputenc}
  \usepackage[T1]{fontenc}
  \usepackage{lmodern}
\else
  \usepackage{fontspec}   % Latin Modern (Computer Modern OpenType) por defecto
\fi
\usepackage[spanish,es-tabla,es-nodecimaldot,es-noshorthands,es-lcroman]{babel}

% --- TIPOGRAFÍA (Latin Modern = Computer Modern vectorial) ----------------
\usepackage[final]{microtype}
\linespread{1.06}
\emergencystretch=2em

% --- MATEMÁTICAS -------------------------------------------------------------
\usepackage{amsmath, amssymb, mathtools}
\usepackage{bm}
\usepackage{mathrsfs}
\usepackage{cancel}

% --- PÁGINA -----------------------------------------------------------------
\usepackage[a4paper, top=2.9cm, bottom=2.7cm, left=2.7cm, right=2.7cm,
            headsep=0.9cm, headheight=22pt]{geometry}
\usepackage{graphicx}

% --- PALETA OSCURA ------------------------------------------------------------
\usepackage[table]{xcolor}
\definecolor{fondo}{HTML}{000000}      % fondo de página (negro puro)
\definecolor{texto}{HTML}{ECECEC}      % texto principal (blanco suave)
\definecolor{blanco}{HTML}{FFFFFF}
\definecolor{negro}{HTML}{000000}
\definecolor{gris}{HTML}{9CA0A8}       % texto secundario, ejes
\definecolor{grismedio}{HTML}{5B5F67}  % reglas, marcos tenues
\definecolor{grisoscuro}{HTML}{2E3036} % rejillas, divisiones
\definecolor{cajafondo}{HTML}{0E1014}  % interior de las cajas
\definecolor{cian}{HTML}{20B4CE}       % definiciones y acento principal
\definecolor{cianclaro}{HTML}{7FDBEA}  % enlaces
\definecolor{violeta}{HTML}{A48BFA}    % teoremas, proposiciones, lemas
\definecolor{ambar}{HTML}{F4B740}      % resultados clave, receta
\definecolor{magenta}{HTML}{E5579D}
\definecolor{verde}{HTML}{4FD68A}
\definecolor{naranja}{HTML}{FA8F45}
\definecolor{rojo}{HTML}{F25F5C}

% Página negra y texto claro. Se fija en el preámbulo para que el color
% «por defecto» (el que restauran flotantes, notas, cabeceras...) sea el claro.
\pagecolor{fondo}
\color{texto}

% --- GRÁFICOS -----------------------------------------------------------------
\usepackage{tikz}
\usetikzlibrary{arrows.meta, calc, positioning, decorations.pathmorphing,
                decorations.markings, decorations.pathreplacing, patterns,
                shapes.geometric, shapes.misc, fit, backgrounds, shadings,
                fadings, intersections}
\usepackage{pgfplots}
\pgfplotsset{compat=1.18}
\usepgfplotslibrary{colormaps, fillbetween, groupplots}

% --- LISTAS, TABLAS, CABECERAS ----------------------------------------------
\usepackage{booktabs}
\usepackage{array}
\usepackage{multirow}
\usepackage{enumitem}
\setlist[itemize]{leftmargin=1.6em, itemsep=2pt, topsep=4pt,
                  label={\color{cian}\scriptsize$\blacktriangleright$}}
\setlist[enumerate]{leftmargin=1.9em, itemsep=2pt, topsep=4pt}
\arrayrulecolor{grismedio}

\usepackage{fancyhdr}
\pagestyle{fancy}
\fancyhf{}
\fancyhead[LE]{\small\color{gris}\thepage\hfill\scshape Ecuaciones Diferenciales}
\fancyhead[RO]{\small\color{gris}\scshape Tema 5\,·\,Puntos ordinarios y series enteras\hfill\upshape\thepage}
\fancyhead[LO,RE]{}
\renewcommand{\headrulewidth}{0.3pt}
\renewcommand{\headrule}{\hbox to\headwidth{\color{grisoscuro}\leaders\hrule height \headrulewidth\hfill}}
\fancypagestyle{plain}{\fancyhf{}\renewcommand{\headrulewidth}{0pt}%
  \fancyfoot[C]{\small\color{gris}\thepage}}

% --- TÍTULOS DE SECCIÓN -----------------------------------------------------
\usepackage{titlesec}
\titleformat{\section}[block]
  {\normalfont\Large\bfseries\color{blanco}}
  {\color{cian}\thesection}{0.9em}{}
  [{\vspace{3pt}\color{grisoscuro}\titlerule[0.5pt]}]
\titlespacing*{\section}{0pt}{4.2ex plus 1ex minus .4ex}{2.2ex plus .4ex}
\titleformat{\subsection}[block]
  {\normalfont\large\bfseries\color{blanco}}
  {\color{cian}\thesubsection}{0.8em}{}
\titlespacing*{\subsection}{0pt}{3.2ex plus 1ex minus .4ex}{1.4ex plus .3ex}
\titleformat{\paragraph}[runin]
  {\normalfont\bfseries\color{blanco}}{}{0pt}{}[.]
\titlespacing*{\paragraph}{0pt}{1.6ex plus .5ex}{0.6em}
% Evita títulos huérfanos al pie de página
\usepackage{needspace}
\usepackage{placeins}
\usepackage{etoolbox}
\AtBeginDocument{%
  \pretocmd{\section}{\needspace{9\baselineskip}}{}{}%
  \pretocmd{\subsection}{\needspace{6\baselineskip}}{}{}}
\setcounter{secnumdepth}{2}
\setcounter{tocdepth}{2}

% --- ÍNDICE -------------------------------------------------------------------
\usepackage{tocloft}
\renewcommand{\cfttoctitlefont}{\normalfont\Large\bfseries\color{blanco}}
\renewcommand{\cftsecfont}{\bfseries\color{blanco}}
\renewcommand{\cftsecpagefont}{\bfseries\color{cian}}
\renewcommand{\cftsubsecfont}{\color{texto}}
\renewcommand{\cftsubsecpagefont}{\color{gris}}
\renewcommand{\cftsecleader}{\color{grismedio}\cftdotfill{\cftdotsep}}
\renewcommand{\cftsubsecleader}{\color{grisoscuro}\cftdotfill{\cftdotsep}}
\setlength{\cftbeforesecskip}{6pt}
\renewcommand{\contentsname}{Índice}

% --- ENLACES Y REFERENCIAS CRUZADAS ----------------------------------------
\usepackage[colorlinks=true, linkcolor=cianclaro, citecolor=ambar,
            urlcolor=cianclaro, plainpages=false, pdfusetitle=false,
            pdftitle={Tema 5. Puntos ordinarios y series enteras},
            pdfauthor={Daniel Martín Guerrero},
            pdfsubject={Ecuaciones Diferenciales},
            bookmarksnumbered=true]{hyperref}
\usepackage{amsthm}
\usepackage[spanish,nameinlink]{cleveref}   % tras amsthm y hyperref

% --- BIBLIOGRAFÍA -----------------------------------------------------------
\usepackage{csquotes}
\usepackage[backend=biber, style=numeric, sorting=none, giveninits=true,
            maxbibnames=6, minbibnames=3, doi=false, url=false,
            isbn=false]{biblatex}
\addbibresource{referencias.bib}
\renewcommand*{\bibfont}{\small}
\setlength{\bibitemsep}{4pt}

% --- PIES DE FIGURA -------------------------------------------------------
\usepackage[font={small,color=texto}, labelfont={sc,color=cian},
            labelsep=period, margin=0.4cm, justification=justified,
            skip=8pt]{caption}
\usepackage{subcaption}
\numberwithin{equation}{section}
\numberwithin{figure}{section}
\numberwithin{table}{section}

% =========================================================================
%               NOTACIÓN Y MACROS
% =========================================================================
\newcommand{\dif}{\mathrm{d}}
\newcommand{\eul}{\mathrm{e}}
\newcommand{\iu}{\mathrm{i}}
\newcommand{\R}{\mathbb{R}}
\newcommand{\C}{\mathbb{C}}
\newcommand{\N}{\mathbb{N}}
\newcommand{\Z}{\mathbb{Z}}
\newcommand{\abs}[1]{\lvert #1\rvert}
\newcommand{\Abs}[1]{\left\lvert #1\right\rvert}
\newcommand{\defeq}{\mathrel{\overset{\mathrm{def}}{=}}}
\newcommand{\Cinf}{\mathcal{C}^{\infty}}
\newcommand{\Com}{\mathcal{C}^{\omega}}           % funciones analíticas reales
\newcommand{\disco}[2]{D(#1,#2)}                  % disco abierto
\newcommand{\dd}[2]{\frac{\dif #1}{\dif #2}}
\newcommand{\ddd}[2]{\frac{\dif^{2} #1}{\dif #2^{2}}}
\newcommand{\serie}[3][n]{\sum_{#1=#2}^{\infty}#3}
\DeclareMathOperator{\sen}{sen}
\DeclareMathOperator{\senh}{senh}
\DeclareMathOperator{\artanh}{artanh}
\DeclareMathOperator{\dist}{dist}
\DeclareMathOperator{\Real}{Re}
\DeclareMathOperator{\Imag}{Im}
\DeclareMathOperator{\Arg}{arg}
\DeclareMathOperator{\ord}{ord}
\DeclareMathOperator{\Ai}{Ai}
\DeclareMathOperator{\Bi}{Bi}
\renewcommand{\sin}{\sen}
\newcommand{\Heav}{\Theta}
\newcommand{\nota}[1]{\textcolor{gris}{#1}}        % texto secundario
\newcommand{\resalta}[1]{\textcolor{blanco}{\bfseries #1}}
% Llaves horizontales en color tenue
\newcommand{\ubr}[2]{\underbrace{#1}_{\textcolor{gris}{#2}}}

% =========================================================================
%        ENTORNOS FORMALES
% =========================================================================
\usepackage[most]{tcolorbox}

% Contador común a todos los enunciados (numeración «sección.n»)
\newcounter{enunciado}[section]
\renewcommand{\theenunciado}{\thesection.\arabic{enunciado}}

% --- Cajas con pestaña (definiciones, teoremas, resultados) -------------------
% #1 color del marco, #2 rótulo completo
\tcbset{
  caja/.style={enhanced, colback=cajafondo, colupper=texto,
    colframe=#1, boxrule=0.7pt, arc=2.5pt, outer arc=2.5pt,
    left=10pt, right=10pt, top=11pt, bottom=8pt,
    before skip=16pt plus 2pt, after skip=14pt plus 2pt,
    attach boxed title to top left={yshift*=-\tcboxedtitleheight/2, xshift=10pt},
    boxed title style={colback=#1, colframe=#1, arc=2pt, outer arc=2pt,
                       boxrule=0pt, left=5pt, right=5pt, top=1.5pt, bottom=1.5pt},
    fonttitle=\sffamily\bfseries\small, coltitle=negro},
}
% Constructor: \cajaenunciado{nombre}{Rótulo}{color}{estilo cuerpo}
\makeatletter
\newcommand{\cajaenunciado}[4]{%
  \newcounter{#1}%
  \expandafter\def\csname the#1\endcsname{\theenunciado}%
  \expandafter\def\csname theH#1\endcsname{#1.\arabic{#1}}%
  \newtcolorbox{tcb@#1}[1]{caja=#3,
      title={#2~\theenunciado\ifx\relax##1\relax\else\ \textnormal{\sffamily(##1)}\fi}}%
  \newenvironment{#1}[1][]%
    {\par\stepcounter{enunciado}\refstepcounter{#1}%
     \begin{tcb@#1}{##1}#4\ignorespaces}%
    {\end{tcb@#1}}%
}
\makeatother
\cajaenunciado{definicion}{Definición}{cian}{\normalfont}
\cajaenunciado{teorema}{Teorema}{violeta}{\itshape}
\cajaenunciado{proposicion}{Proposición}{violeta}{\itshape}
\cajaenunciado{lema}{Lema}{violeta}{\itshape}
\cajaenunciado{corolario}{Corolario}{violeta}{\itshape}
\cajaenunciado{resultado}{Resultado}{ambar}{\normalfont}

% Caja de método (receta), sin numerar
\newtcolorbox{metodo}[1][Método]{caja=ambar, title={#1}}

% --- Enunciados sin caja (ejemplos, observaciones, notación) -----------------
\newtheoremstyle{ejemplo}{12pt}{10pt}{\normalfont}{}%
  {\bfseries\color{cian}}{.}{0.6em}%
  {\thmname{#1}\thmnumber{ #2}\thmnote{\normalfont\color{gris}\ (#3)}}
\newtheoremstyle{observacion}{10pt}{10pt}{\normalfont}{}%
  {\itshape\bfseries\color{violeta}}{.}{0.6em}%
  {\thmname{#1}\thmnumber{ #2}\thmnote{\normalfont\itshape\color{gris}\ (#3)}}
\theoremstyle{ejemplo}
\newtheorem{ejemplo}[enunciado]{Ejemplo}
\newtheorem{notacion}[enunciado]{Notación}
\theoremstyle{observacion}
\newtheorem{observacion}[enunciado]{Observación}

% Demostraciones
\renewcommand{\qedsymbol}{\textcolor{gris}{$\blacksquare$}}
\makeatletter
\renewenvironment{proof}[1][\proofname]{\par
  \pushQED{\qed}\normalfont\topsep6\p@\@plus6\p@\relax
  \trivlist\item[\hskip\labelsep\itshape\bfseries\color{gris}#1\@addpunct{.}]\ignorespaces
}{\popQED\endtrivlist\@endpefalse}
\makeatother
\newenvironment{esbozo}{\begin{proof}[Esbozo de la demostración]}{\end{proof}}

% Notas al margen del discurso: «idea», recuadro tenue con barra lateral
\newtcolorbox{idea}[1][Idea]{enhanced, breakable, colback=cajafondo, colupper=texto,
  frame hidden, borderline west={2pt}{0pt}{cian!80}, arc=0pt,
  left=10pt, right=8pt, top=6pt, bottom=6pt, before skip=12pt, after skip=12pt,
  fontupper=\normalfont,
  before upper={{\sffamily\bfseries\small\color{cian}#1.}\ }}

\crefname{definicion}{definición}{definiciones}
\Crefname{definicion}{Definición}{Definiciones}
\crefname{teorema}{teorema}{teoremas}
\Crefname{teorema}{Teorema}{Teoremas}
\crefname{proposicion}{proposición}{proposiciones}
\Crefname{proposicion}{Proposición}{Proposiciones}
\crefname{lema}{lema}{lemas}
\Crefname{lema}{Lema}{Lemas}
\crefname{corolario}{corolario}{corolarios}
\Crefname{corolario}{Corolario}{Corolarios}
\crefname{resultado}{resultado}{resultados}
\Crefname{resultado}{Resultado}{Resultados}
\crefname{ejemplo}{ejemplo}{ejemplos}
\Crefname{ejemplo}{Ejemplo}{Ejemplos}
\crefname{observacion}{observación}{observaciones}
\Crefname{observacion}{Observación}{Observaciones}
\crefname{notacion}{notación}{notaciones}
\Crefname{notacion}{Notación}{Notaciones}
\crefname{section}{sección}{secciones}
\Crefname{section}{Sección}{Secciones}
\crefname{subsection}{apartado}{apartados}
\Crefname{subsection}{Apartado}{Apartados}
\crefname{figure}{figura}{figuras}
\Crefname{figure}{Figura}{Figuras}
\crefname{table}{tabla}{tablas}
\Crefname{table}{Tabla}{Tablas}
\crefname{equation}{ecuación}{ecuaciones}
\Crefname{equation}{Ecuación}{Ecuaciones}
\crefformat{equation}{(#2#1#3)}
\Crefformat{equation}{Ecuación~(#2#1#3)}
\crefrangeformat{equation}{(#3#1#4)--(#5#2#6)}
\crefmultiformat{equation}{(#2#1#3)}{ y~(#2#1#3)}{, (#2#1#3)}{ y~(#2#1#3)}

% Modo rápido: si existe el fichero vacío «rapido.flag», las superficies y
% retratos de fase se generan con baja resolución (compilación en segundos).
\newif\ifrapido
\IfFileExists{rapido.flag}{\rapidotrue}{\rapidofalse}
\ifrapido
  \def\resPx{24}\def\resPy{34}\def\resF{41}
\else
  \def\resPx{85}\def\resPy{131}\def\resF{141}
\fi

% =========================================================================
%        ESTILOS GRÁFICOS COMUNES
% =========================================================================
% =========================================================================
%  Estilos gráficos comunes (TikZ / pgfplots) para fondo negro
% =========================================================================

% --- Mapa de colores «arcoíris» para coloreado por fase (dominio complejo) ---
% Recorre el círculo cromático completo: el color codifica arg f(z) en (-pi,pi].
\pgfplotsset{colormap={fase}{
  rgb255=(255,40,60) rgb255=(255,150,20) rgb255=(250,230,40)
  rgb255=(60,220,90) rgb255=(30,200,230) rgb255=(60,90,255)
  rgb255=(190,60,240) rgb255=(255,40,60)}}

% --- Estilos TikZ generales ------------------------------------------------
\tikzset{
  >=Stealth,
  every picture/.append style={line cap=round, line join=round},
  eje/.style={-{Stealth[length=5pt,width=4pt]}, line width=0.6pt, gris},
  ejeclaro/.style={-{Stealth[length=5pt,width=4pt]}, line width=0.7pt, texto},
  fino/.style={line width=0.4pt, grismedio},
  guia/.style={densely dashed, line width=0.45pt, grismedio},
  punteado/.style={densely dotted, line width=0.6pt, gris},
  etiq/.style={font=\small, text=texto},
  etiqgris/.style={font=\footnotesize, text=gris},
  polo/.style={cross out, draw=rojo, line width=1.1pt, minimum size=6pt, inner sep=0pt},
  punto/.style={circle, fill=blanco, inner sep=0pt, minimum size=4pt},
  puntocian/.style={circle, fill=cian, inner sep=0pt, minimum size=4.5pt},
  orienta/.style={decoration={markings, mark=at position #1 with
     {\arrow{Stealth[length=6pt,width=5pt]}}}, postaction={decorate}},
  orienta/.default=0.55,
  nodo/.style={draw=grismedio, fill=cajafondo, rounded corners=2pt,
     line width=0.5pt, inner xsep=6pt, inner ysep=4pt, font=\small, text=texto},
  nodocian/.style={nodo, draw=cian},
  nodoambar/.style={nodo, draw=ambar},
  nodovioleta/.style={nodo, draw=violeta},
  flecha/.style={-{Stealth[length=5pt,width=4.5pt]}, line width=0.6pt, gris},
}

% --- Estilo base de ejes pgfplots en oscuro ----------------------------------
\pgfplotsset{
  oscuro/.style={
    axis background/.style={fill=fondo},
    tick label style={font=\scriptsize, text=gris},
    label style={font=\small, text=texto},
    title style={font=\small, text=texto},
    axis line style={grismedio, line width=0.5pt},
    every tick/.style={grismedio},
    major tick length=3pt,
    grid=major, grid style={grisoscuro, line width=0.35pt},
    legend style={font=\scriptsize, draw=grisoscuro, fill=cajafondo,
                  text=texto, rounded corners=1.5pt, inner sep=3pt,
                  row sep=1pt},
    legend cell align=left,
    /pgf/number format/use comma,
    /pgf/number format/1000 sep={},
    unbounded coords=jump,
  },
  % Ejes cruzados por el origen (estilo «pizarra»)
  cruz/.style={oscuro, axis lines=middle, grid=none,
    axis line style={gris, line width=0.6pt,
      -{Stealth[length=5pt,width=4pt]}},
    xlabel style={anchor=west, font=\small},
    ylabel style={anchor=south, font=\small},
  },
  superficie/.style={
    tick label style={font=\scriptsize, text=gris},
    label style={font=\small, text=texto},
    axis line style={grismedio}, every tick/.style={grismedio},
    grid=none, colormap name=fase, shader=interp,
    /pgf/number format/use comma,
  },
  % Curvas
  c1/.style={line width=1.15pt, cian},
  c2/.style={line width=1.15pt, magenta},
  c3/.style={line width=1.15pt, ambar},
  c4/.style={line width=1.15pt, verde},
  c5/.style={line width=1.15pt, violeta},
  c6/.style={line width=1.15pt, naranja},
}




% =========================================================================
%                     VARIABLES PARA LA PORTADA
% =========================================================================
\newcommand{\Universidad}{Universidad de Salamanca}
\newcommand{\Facultad}{Facultad de Ciencias}
\newcommand{\Titulacion}{Grado en Física}
\newcommand{\Asignatura}{Ecuaciones Diferenciales}
\newcommand{\NumTema}{Tema 5}
\newcommand{\TituloTema}{EDO lineales de segundo orden\\con coeficientes variables}
\newcommand{\SubtituloTema}{Parte I\,·\,Puntos ordinarios y soluciones en serie entera}
\newcommand{\tipotrabajo}{Apuntes formalizados}
\newcommand{\miNombre}{Daniel Martín Guerrero}
\newcommand{\miProfesorado}{\begin{tabular}{@{}c@{}}Pilar García Estévez\\José Beltrán Jiménez\\María Pérez Garrote\end{tabular}}
\newcommand{\miUbicacion}{Salamanca, curso 2026/2027}

% =========================================================================
%                                DOCUMENTO
% =========================================================================
\begin{document}

% --- PORTADA ----------------------------------------------------------------
\pagenumbering{alph}
%%%%%%%%%%%%%%%%%%%%%%%%%%%%%%%%%%%%%%%%%%%%%%%%%%%%%%%%%%%%%%%%%%%%%%%%
% Portada oscura: identificación institucional, título, superficie
% |1/(1+z^2)| coloreada por la fase (los dos polos z = ±i que limitan el
% radio de convergencia de la serie real) y datos de autoría.
%%%%%%%%%%%%%%%%%%%%%%%%%%%%%%%%%%%%%%%%%%%%%%%%%%%%%%%%%%%%%%%%%%%%%%%%
\begin{titlepage}
\thispagestyle{empty}
\begin{tikzpicture}[remember picture, overlay]
  \coordinate (O) at (current page.south west);
  \coordinate (C) at ($(O)+(10.5cm,0)$);

  % Logotipo y cabecera institucional
  \node[anchor=north, inner sep=0] at ($(C)+(0,27.6cm)$)
    {\includegraphics[height=1.45cm]{logo_usal_blanco.png}};
  \draw[grismedio, line width=0.35pt] ($(O)+(3.0cm,25.6cm)$) -- ($(O)+(18.0cm,25.6cm)$);
  \node[anchor=north, align=center, inner sep=0] at ($(C)+(0,25.15cm)$)
    {{\color{gris}\scshape\small\textls[140]{\Facultad\ \ ·\ \ \Titulacion}}};

  % Asignatura, tema y título
  \node[anchor=north, align=center, inner sep=0, text width=16cm] at ($(C)+(0,23.3cm)$)
    {{\color{cian}\scshape\large\textls[180]{\Asignatura}}\\[6pt]
     {\color{gris}\scshape\normalsize\textls[140]{\NumTema}}\\[20pt]
     {\color{blanco}\fontsize{27}{34}\selectfont\bfseries\TituloTema}\\[14pt]
     {\color{cian}\rule{3.2cm}{0.7pt}}\\[12pt]
     {\color{texto}\itshape\large\SubtituloTema}};

  % Superficie |f(z)| de f(z) = 1/(1+z^2), coloreada por arg f(z)
  \node[anchor=center, inner sep=0] at ($(C)+(0,10.9cm)$) {%
    \begin{tikzpicture}
      % La fase salta de +180 a -180 grados sobre el eje imaginario |y|>1:
      % se dibuja la superficie en dos mitades (primero la lejana, x>0) para que el
      % sombreado interpolado no recorra todo el círculo cromático en el corte.
      \begin{axis}[superficie, width=15.5cm, height=11.5cm,
          view={-28}{24}, axis lines=none, clip=false,
          y domain=-2:2, samples=\resPx, samples y=\resPy,
          zmin=0, zmax=4.6, point meta min=-180, point meta max=180]
        \addplot3[surf, z buffer=sort, domain=0.0001:2,
            point meta={-atan2(2*x*y,1+x^2-y^2)}]
          {min(1/sqrt((1+x^2-y^2)^2+4*x^2*y^2+0.00005),4.6)};
        \addplot3[surf, z buffer=sort, domain=-2:-0.0001,
            point meta={-atan2(2*x*y,1+x^2-y^2)}]
          {min(1/sqrt((1+x^2-y^2)^2+4*x^2*y^2+0.00005),4.6)};
      \end{axis}
    \end{tikzpicture}};

  % Epígrafe: la serie entera y la cota del radio
  \node[anchor=north, align=center, inner sep=0] at ($(C)+(0,5.5cm)$)
    {\color{texto}$\displaystyle
      y(x)=\sum_{n=0}^{\infty}a_n\,(x-x_0)^n,\qquad
      \frac{1}{1+x^{2}}=\sum_{n=0}^{\infty}(-1)^n x^{2n}\quad(\abs{x}<1)$};

  % Autoría, lugar y curso
  \draw[grismedio, line width=0.35pt] ($(O)+(3.0cm,4.2cm)$) -- ($(O)+(18.0cm,4.2cm)$);
  \node[anchor=north west, inner sep=0, align=left] at ($(O)+(3.0cm,3.75cm)$)
    {{\scshape\small\color{cian}\textls[100]{Autor}}\\[3pt]
     {\large\color{blanco}\miNombre}};
  \node[anchor=north, inner sep=0, align=center] at ($(C)+(0,3.75cm)$)
    {{\scshape\small\color{cian}\textls[100]{Profesorado}}\\[3pt]
     {\small\color{texto}\miProfesorado}};
  \node[anchor=north east, inner sep=0, align=right] at ($(O)+(18.0cm,3.75cm)$)
    {{\scshape\small\color{cian}\textls[100]{\tipotrabajo}}\\[3pt]
     {\color{gris}\miUbicacion}};
\end{tikzpicture}
\end{titlepage}

\newpage\thispagestyle{empty}\null\newpage
\pagenumbering{arabic}
\setcounter{page}{1}

% --- CABECERA DEL ARTÍCULO ------------------------------------------------
\thispagestyle{plain}
\begin{center}
  {\color{gris}\scshape\small\textls[120]{\Asignatura\ \,·\ \,\NumTema}}\par\vspace{10pt}
  {\color{blanco}\LARGE\bfseries Puntos ordinarios y soluciones\\[3pt]
   en serie entera de EDO lineales de segundo orden}\par\vspace{12pt}
  {\color{cian}\rule{2.8cm}{0.6pt}}\par\vspace{10pt}
  {\large\color{texto}\miNombre}\par\vspace{2pt}
  {\small\itshape\color{gris}\Titulacion, \Facultad, \Universidad}\par
\end{center}
\vspace{14pt}

\begin{center}
\begin{minipage}{0.86\textwidth}
\small
{\scshape\color{cian}Resumen.}\ \
Estudiamos la ecuación lineal homogénea de segundo orden
$P(x)\,y''+A(x)\,y'+B(x)\,y=0$ con coeficientes variables en un entorno de
un \emph{punto ordinario} $x_0$, esto es, un punto en el que $P(x_0)\neq 0$.
Tras motivar el método a través de la analogía entre ecuaciones diferenciales
y ecuaciones en diferencias, fijamos los fundamentos analíticos
---funciones analíticas, derivación término a término, unicidad del
desarrollo en serie y clasificación de singularidades aisladas, con especial
atención a los polos--- y explicamos por qué las singularidades
\emph{complejas} de los coeficientes acotan el radio de convergencia de las
soluciones \emph{reales}. Enunciamos y demostramos el teorema de existencia
de soluciones analíticas (teorema de Fuchs para puntos ordinarios) mediante
el método de las series mayorantes, deducimos la relación de recurrencia
general para los coeficientes y construimos un sistema fundamental de
soluciones. El método se ilustra con el oscilador armónico, la ecuación de
Airy, un problema de valores iniciales con solución gaussiana y la ecuación
de Legendre, cuyo truncamiento conduce a los polinomios de Legendre.

\medskip
{\scshape\color{cian}Palabras clave.}\ \
función analítica, polo, serie de Laurent, radio de convergencia, punto
ordinario, teorema de Fuchs, relación de recurrencia, producto de Cauchy,
wronskiano, ecuación de Airy, ecuación de Legendre.

\medskip
{\scshape\color{cian}MSC 2020.}\ \ 34A25, 34A30, 34M05, 30B10, 30D10.
\end{minipage}
\end{center}

\clearpage
{\hypersetup{linkcolor=texto}\tableofcontents}
\clearpage

% =========================================================================
\section{Motivación: la analogía con las ecuaciones en diferencias}\label{sec:motivacion}
% =========================================================================

Las ecuaciones diferenciales lineales con coeficientes \emph{constantes} se
resuelven con el ansatz exponencial $y=\eul^{\lambda x}$, que transforma la
ecuación en un polinomio característico. Existe un mecanismo idéntico para
las ecuaciones \emph{discretas}, y es precisamente ese mundo discreto el que
reaparecerá, de forma inesperada, cuando busquemos soluciones en serie de
ecuaciones con coeficientes \emph{variables}.

% -------------------------------------------------------------------------
\subsection{Ecuaciones en diferencias con coeficientes constantes}\label{sec:diferencias}
% -------------------------------------------------------------------------

\begin{definicion}[Ecuación en diferencias lineal]\label{def:diferencias}
Sean $k\in\N$, $k\ge 1$, y $a_0,\dots,a_k\in\C$ con $a_k\neq 0$ y $a_0\neq0$.
Una \emph{ecuación en diferencias lineal homogénea de orden $k$ con
coeficientes constantes} es una relación
\begin{equation}\label{eq:diferencias}
  a_k\,f(n+k)+a_{k-1}\,f(n+k-1)+\dots+a_1\,f(n+1)+a_0\,f(n)=0,
  \qquad n\in\N ,
\end{equation}
cuya incógnita es una sucesión $f\colon\N\to\C$. Su \emph{polinomio
característico} es $\chi(\lambda)=a_k\lambda^k+a_{k-1}\lambda^{k-1}+\dots+a_0$.
\end{definicion}

Como $a_k\neq 0$, la ecuación \eqref{eq:diferencias} permite despejar
$f(n+k)$ en función de los $k$ valores anteriores: fijados
$f(0),\dots,f(k-1)$, la sucesión queda determinada por inducción. El conjunto
de soluciones es, por tanto, un espacio vectorial de dimensión $k$, en
perfecto paralelismo con el teorema de existencia y unicidad de las EDO
lineales.

\begin{proposicion}[Ansatz exponencial discreto]\label{prop:ansatz-discreto}
La sucesión $f(n)=\lambda^n$, con $\lambda\neq0$, es solución de
\eqref{eq:diferencias} si y solo si $\chi(\lambda)=0$. Si $\chi$ tiene $k$
raíces distintas $\lambda_1,\dots,\lambda_k$, la solución general es
\begin{equation}\label{eq:sol-diferencias}
  f(n)=C_1\lambda_1^{\,n}+C_2\lambda_2^{\,n}+\dots+C_k\lambda_k^{\,n},
  \qquad C_1,\dots,C_k\in\C .
\end{equation}
Si $\lambda_j$ es raíz de multiplicidad $m_j$, las sucesiones
$n^{i}\lambda_j^{\,n}$, $0\le i<m_j$, sustituyen a las potencias repetidas.
\end{proposicion}

\begin{proof}
Sustituyendo $f(n)=\lambda^n$ y extrayendo el factor común $\lambda^n\neq0$,
\begin{equation}\label{eq:caracteristico}
  \lambda^n\bigl(a_k\lambda^k+a_{k-1}\lambda^{k-1}+\dots+a_0\bigr)=\lambda^n\chi(\lambda)=0
  \iff \chi(\lambda)=0 .
\end{equation}
Por linealidad, \eqref{eq:sol-diferencias} es solución. Para ver que es la
general basta comprobar que se pueden imponer datos iniciales arbitrarios
$f(0),\dots,f(k-1)$: el sistema lineal resultante para $(C_1,\dots,C_k)$
tiene por matriz la de Vandermonde $\bigl(\lambda_j^{\,i}\bigr)_{0\le i<k,\,1\le j\le k}$,
cuyo determinante $\prod_{i<j}(\lambda_j-\lambda_i)$ no se anula por ser las
raíces distintas. El caso de raíces múltiples es análogo al de las EDO y
se omite (véase \cite[cap.~2]{ELAYDI}).
\end{proof}

\begin{ejemplo}[Sucesión de Fibonacci]\label{ej:fibonacci}
La ecuación $f(n+2)=f(n+1)+f(n)$, con $f(0)=0$, $f(1)=1$, tiene polinomio
característico $\lambda^2-\lambda-1$, de raíces
$\varphi=\tfrac{1+\sqrt5}{2}$ y $\psi=\tfrac{1-\sqrt5}{2}$. Imponiendo los
datos iniciales en $f(n)=C_1\varphi^n+C_2\psi^n$ se obtiene la fórmula de
Binet,
\[
  f(n)=\frac{\varphi^{n}-\psi^{n}}{\sqrt5},
\]
exactamente como se resolvería $y''-y'-y=0$ con $y=\eul^{\lambda x}$.
\end{ejemplo}

El paralelismo es completo y se resume en la \cref{tab:analogia}: el
operador de derivación $D=\dif/\dif x$ desempeña en el continuo el papel del
operador de desplazamiento $E\colon f(n)\mapsto f(n+1)$ en el discreto.

\begin{table}[ht]
\centering
\caption{Diccionario entre las EDO lineales con coeficientes constantes y
las ecuaciones en diferencias lineales con coeficientes constantes.}
\label{tab:analogia}
\renewcommand{\arraystretch}{1.35}
\begin{tabular}{@{}l>{\centering\arraybackslash}p{4.6cm}>{\centering\arraybackslash}p{4.6cm}@{}}
\toprule
 & \textcolor{cian}{\bfseries Continuo (EDO)} & \textcolor{cian}{\bfseries Discreto (diferencias)}\\
\midrule
Incógnita            & $y(x)$, $x\in\R$                     & $f(n)$, $n\in\N$\\
Operador básico      & $D\colon y\mapsto y'$                & $E\colon f(n)\mapsto f(n+1)$\\
Ecuación             & $\sum_{j} a_j\,D^{j}y=0$             & $\sum_{j} a_j\,E^{j}f=0$\\
Ansatz               & $y=\eul^{\lambda x}$                 & $f(n)=\lambda^{n}$\\
Polinomio característico & $\sum_j a_j\lambda^j=0$          & $\sum_j a_j\lambda^j=0$\\
Raíz de multiplicidad $m$ & $x^{i}\eul^{\lambda x}$, $i<m$  & $n^{i}\lambda^{n}$, $i<m$\\
Datos que fijan la solución & $y(x_0),\dots,y^{(k-1)}(x_0)$ & $f(0),\dots,f(k-1)$\\
\bottomrule
\end{tabular}
\end{table}

\begin{observacion}[$E=\eul^{D}$]\label{obs:E-expD}
El diccionario no es una mera coincidencia formal. Si $y$ es una función
cuya serie de Taylor en $x$ converge a $y$ en un radio mayor que $1$,
\[
  y(x+1)=\sum_{k=0}^{\infty}\frac{y^{(k)}(x)}{k!}
        =\Bigl(\sum_{k=0}^{\infty}\frac{D^{k}}{k!}\Bigr)y(x)
        \eqqcolon \eul^{D}y(x),
\]
de modo que $E=\eul^{D}$ y el ansatz $\lambda^n=\eul^{n\ln\lambda}$ es el
ansatz exponencial evaluado en los enteros. Esta identidad exige que $y$ sea
\emph{analítica} (\cref{def:analitica}): la noción central de este tema
aparece ya en la motivación.
\end{observacion}

% -------------------------------------------------------------------------
\subsection{Por qué nos interesa: de la EDO a la recurrencia}\label{sec:por-que}
% -------------------------------------------------------------------------

Cuando los coeficientes de la EDO dependen de $x$, el ansatz
$y=\eul^{\lambda x}$ deja de funcionar. La estrategia de este tema consiste
en proponer una \emph{serie de potencias} $y=\sum a_n(x-x_0)^n$ y traducir la
ecuación diferencial para $y(x)$ en una ecuación en diferencias ---una
\emph{relación de recurrencia}--- para la sucesión de coeficientes
$(a_n)_{n\ge0}$. Cambiamos así un problema continuo por uno discreto
(\cref{fig:analogia}). Dos diferencias con respecto a la
\cref{def:diferencias} merecen destacarse desde ahora:
\begin{enumerate}[label=(\roman*)]
  \item los coeficientes de la recurrencia \emph{dependen de $n$}; por
        ejemplo, $y''+y=0$ conduce a $(n+2)(n+1)\,a_{n+2}+a_n=0$;
  \item si los coeficientes de la EDO no son polinómicos, la recurrencia
        involucra \emph{todos} los coeficientes anteriores (es de orden no
        acotado), como veremos en la \cref{sec:recurrencia}.
\end{enumerate}

% Figura: esquema del método (continuo -> discreto -> continuo)
\begin{figure}[ht]
\centering
\begin{tikzpicture}
  \node[nodocian, align=center, text width=4.1cm, minimum height=1.5cm] (edo) at (0,0)
    {\textcolor{cian}{\bfseries EDO}\\[3pt] $P\,y''+A\,y'+B\,y=0$\\[2pt]
     \footnotesize\color{gris} incógnita: la función $y(x)$};
  \node[nodoambar, align=center, text width=4.6cm, minimum height=1.5cm] (rec) at (8.6,0)
    {\textcolor{ambar}{\bfseries Recurrencia}\\[3pt] $a_{n+2}=\Phi_n(a_0,\dots,a_{n+1})$\\[2pt]
     \footnotesize\color{gris} incógnita: la sucesión $(a_n)_{n\ge0}$};
  \node[nodovioleta, align=center, minimum height=1.1cm, inner xsep=10pt] (sol) at (4.3,-2.9)
    {\textcolor{violeta}{\bfseries Solución}\quad
     $y(x)=\sum_{n\ge0}a_n(x-x_0)^n$};

  \draw[flecha, cian] (edo.east) -- (rec.west)
     node[midway, above=3pt, etiqgris] {ansatz en serie}
     node[midway, below=3pt, etiqgris] {continuo $\to$ discreto};
  \draw[flecha, ambar] (rec.south) |- (sol.east)
     node[pos=0.25, right=4pt, etiqgris, align=left] {resolver\\por inducción};
  \draw[flecha, violeta] (sol.west) -| (edo.south)
     node[pos=0.75, left=4pt, etiqgris, align=right] {sumar la serie\\(radio $R>0$)};
\end{tikzpicture}
\caption{Estructura del método de series. La ecuación diferencial para
$y(x)$ se convierte en una ecuación en diferencias para los coeficientes
$a_n$; resuelta ésta, la serie reconstruye la solución en el disco de
convergencia.}
\label{fig:analogia}
\end{figure}


% =========================================================================
\section{Planteamiento del problema y funciones analíticas}\label{sec:planteamiento}
% =========================================================================

% -------------------------------------------------------------------------
\subsection{La ecuación}\label{sec:ecuacion}
% -------------------------------------------------------------------------
Nos limitamos a ecuaciones lineales homogéneas de segundo orden que, tras
«quitar denominadores», adoptan la forma
\begin{equation}\label{eq:edo}
  \boxed{\;P(x)\,y''(x)+A(x)\,y'(x)+B(x)\,y(x)=0\;}
\end{equation}
sobre un intervalo abierto $I\subseteq\R$, donde $P$ es un polinomio de grado
bajo (habitualmente dos, a lo sumo tres) y $A$, $B$ son funciones analíticas
en $I$ \cite{TEMAV}. Son ejemplos típicos
\[
  y''-x^{3}y'+y=0,\qquad
  y''+(\sin x)\,y=0,\qquad
  (1-x^{2})\,y''-2x\,y'+\ell(\ell+1)\,y=0 ,
\]
la última de las cuales es la \emph{ecuación de Legendre}. Para las
ecuaciones completas $P\,y''+A\,y'+B\,y=g(x)$ se aplican los métodos ya
conocidos (variación de las constantes) para obtener una solución
particular, de modo que basta estudiar la homogénea.

Los coeficientes ya no son constantes y el ansatz exponencial fracasa: en
$y''-x\,y=0$, por ejemplo, $y=\eul^{\lambda x}$ conduce a
$(\lambda^{2}-x)\,\eul^{\lambda x}=0$, que no puede cumplirse para todo $x$
con $\lambda$ constante. Necesitamos otra herramienta.

\begin{teorema}[Existencia y unicidad; recordatorio]\label{teo:eyu}
Sea $J\subseteq I$ un intervalo abierto en el que $P$ no se anula, y sea
$x_0\in J$. Para cualesquiera $\alpha,\beta\in\R$ existe una única solución
$y\in\mathcal C^{2}(J)$ de \eqref{eq:edo} con $y(x_0)=\alpha$,
$y'(x_0)=\beta$. En consecuencia, las soluciones de \eqref{eq:edo} en $J$
forman un espacio vectorial real de dimensión $2$
\textup{\cite{ELSGOLTZ}, \cite[cap.~IV]{TEMAV}}.
\end{teorema}

El \cref{teo:eyu} garantiza que las soluciones existen, pero no dice cómo
construirlas. El método de series proporcionará una construcción explícita
y, además, una propiedad cualitativa nueva: la \emph{analiticidad} de las
soluciones.

% -------------------------------------------------------------------------
\subsection{Funciones analíticas}\label{sec:analiticas}
% -------------------------------------------------------------------------

\begin{definicion}[Función analítica]\label{def:analitica}
Una función $f\colon I\to\R$ es \emph{analítica en $x_0\in I$} si existen
$R>0$ y coeficientes $(c_k)_{k\ge0}\subset\R$ tales que
\begin{equation}\label{eq:analitica}
  f(x)=\sum_{k=0}^{\infty}c_k\,(x-x_0)^k
  \qquad\text{para todo } x \text{ con } \abs{x-x_0}<R .
\end{equation}
Es decir: \textup{(i)} $f$ admite un desarrollo en serie de potencias
centrado en $x_0$ y \textup{(ii)} dicha serie \emph{converge a la propia
función} en un entorno de $x_0$. Se dice que $f$ es analítica en $I$ si lo es
en cada punto de $I$, y se escribe $f\in\Com(I)$.
\end{definicion}

La herramienta técnica que legitima todo el método es la siguiente.

\begin{teorema}[Series de potencias]\label{teo:series}
Sea $\sum_{k\ge0}c_k(x-x_0)^k$ una serie de potencias y sea $R\in[0,\infty]$
su \emph{radio de convergencia}, dado por la fórmula de Cauchy--Hadamard
\begin{equation}\label{eq:hadamard}
  \frac1R=\limsup_{k\to\infty}\abs{c_k}^{1/k}.
\end{equation}
\begin{enumerate}[label=\textup{(\roman*)}]
  \item La serie converge absolutamente si $\abs{x-x_0}<R$, y uniformemente
        en cada intervalo cerrado $\abs{x-x_0}\le r<R$; diverge si
        $\abs{x-x_0}>R$.
  \item Su suma $f$ es derivable en $(x_0-R,x_0+R)$ y se deriva
        \emph{término a término}:
        $f'(x)=\sum_{k\ge1}k\,c_k(x-x_0)^{k-1}$; la serie derivada tiene
        \emph{el mismo radio} $R$.
  \item En consecuencia $f\in\Cinf(x_0-R,x_0+R)$ y
        $c_k=f^{(k)}(x_0)/k!$: la serie es la serie de Taylor de su suma.
\end{enumerate}
\end{teorema}

\begin{proof}
(i) es el criterio de la raíz aplicado a $\sum\abs{c_k}\abs{x-x_0}^k$, junto
con el criterio M de Weierstrass. (ii) Como $k^{1/k}\to1$,
$\limsup\abs{k\,c_k}^{1/k}=\limsup\abs{c_k}^{1/k}$, de modo que la serie
derivada tiene radio $R$ y converge uniformemente en cada
$\abs{x-x_0}\le r<R$; el teorema de derivación de sucesiones de funciones
con derivadas uniformemente convergentes permite intercambiar suma y
derivada. (iii) Se obtiene iterando (ii) y evaluando en $x=x_0$.
\end{proof}

\begin{corolario}[Unicidad del desarrollo]\label{cor:unicidad}
Si $\sum_{k\ge0}c_k(x-x_0)^k=0$ para todo $\abs{x-x_0}<R$, con $R>0$, entonces
$c_k=0$ para todo $k$. En particular, si dos series de potencias coinciden en
un entorno de $x_0$, sus coeficientes son iguales uno a uno.
\end{corolario}

\begin{proof}
Por el \cref{teo:series}\,(iii), $c_k=f^{(k)}(x_0)/k!$ con $f\equiv0$.
La segunda afirmación se obtiene aplicando la primera a la diferencia.
\end{proof}

El \cref{cor:unicidad} es el principio de \emph{identificación de
coeficientes} que usaremos para obtener la relación de recurrencia. La
siguiente proposición garantiza que las operaciones algebraicas que haremos
con los coeficientes de la EDO no destruyen la analiticidad.

\begin{proposicion}[Álgebra de funciones analíticas]\label{prop:algebra}
Si $f$ y $g$ son analíticas en $x_0$, también lo son $f+g$, $f\,g$ y, si
$g(x_0)\neq0$, el cociente $f/g$. Los radios de convergencia de $f+g$ y $fg$
son al menos $\min(R_f,R_g)$.
\end{proposicion}

\begin{proof}[Comentario sobre la demostración]
Para la suma es inmediato; para el producto se usa el producto de Cauchy de
series absolutamente convergentes (\cref{lem:cauchy}). El cociente se reduce
a probar que $1/g$ es analítica cuando $g(x_0)\neq0$, lo que resulta
inmediato en variable compleja: $1/g$ es holomorfa en un disco donde $g$ no
se anula y toda función holomorfa en un disco es suma de su serie de Taylor
en él \cite[cap.~4]{AHLFORS}, \cite[§1.1]{KRANTZ}.
\end{proof}

\paragraph{Analítica es más que $\Cinf$}
La condición (ii) de la \cref{def:analitica} no es trivial: existen funciones
infinitamente derivables cuya serie de Taylor converge, pero no a la función.

\begin{proposicion}[Una función $\Cinf$ que no es analítica]\label{prop:suave}
La función
\begin{equation}\label{eq:suave}
  f(x)=\begin{cases}\eul^{-1/x^{2}}, & x\neq0,\\ 0, & x=0,\end{cases}
\end{equation}
es de clase $\Cinf(\R)$ y $f^{(k)}(0)=0$ para todo $k\ge0$. Su serie de
Taylor en $0$ es idénticamente nula y converge en todo $\R$, pero sólo
coincide con $f$ en $x=0$. Por tanto $f$ no es analítica en $0$, y
\[
  \Com(I)\subsetneq\Cinf(I).
\]
\end{proposicion}

\begin{proof}
Por inducción, para $x\neq0$ se tiene $f^{(k)}(x)=Q_k(1/x)\,\eul^{-1/x^{2}}$
con $Q_k$ un polinomio: $Q_0=1$ y
$Q_{k+1}(t)=2t^{3}Q_k(t)-t^{2}Q_k'(t)$. Como
$\lim_{x\to0}x^{-m}\eul^{-1/x^{2}}=\lim_{t\to\infty}t^{m}\eul^{-t^{2}}=0$
para todo $m\in\N$, se sigue que $f^{(k)}(x)\to0$ cuando $x\to0$. Si
$f^{(k)}(0)=0$, entonces
\[
  f^{(k+1)}(0)=\lim_{x\to0}\frac{f^{(k)}(x)-0}{x}
  =\lim_{x\to0}\frac1x\,Q_k(1/x)\,\eul^{-1/x^{2}}=0,
\]
lo que completa la inducción. La serie de Taylor es $\sum 0\cdot x^k\equiv0$,
y $f(x)>0$ para $x\neq0$.
\end{proof}

% Figura: e^{-1/x^2}, función C^infinito no analítica
\begin{figure}[ht]
\centering
\begin{tikzpicture}
  \begin{groupplot}[group style={group size=2 by 1, horizontal sep=1.6cm},
      oscuro, width=0.47\textwidth, height=5.0cm, samples=301]
    \nextgroupplot[xmin=-4, xmax=4, ymin=-0.3, ymax=1.08,
        xtick={-4,-2,0,2,4}, ytick={0,0.5,1},
        xlabel={$x$}, ylabel={$f(x)$},
        title={(a) Vista global}]
      \addplot[c1, domain=-4:4] {exp(-1/(x^2+1e-12))};
      \addplot[c3, densely dashed, domain=-4:4] {0};
      \addplot[guia, domain=-4:4, forget plot] {1};
      \node[font=\scriptsize, text=cian, anchor=west] at (axis cs:2.05,0.55) {$\eul^{-1/x^{2}}$};
      \node[font=\scriptsize, text=ambar, anchor=north] at (axis cs:0,-0.06)
        {serie de Taylor en $0$: $\equiv0$};
    \nextgroupplot[xmin=-0.5, xmax=0.5, ymin=-0.0055, ymax=0.02,
        xtick={-0.5,-0.25,0,0.25,0.5},
        ytick={0,0.01,0.02},
        yticklabel style={/pgf/number format/fixed, /pgf/number format/precision=2},
        scaled y ticks=false,
        xlabel={$x$},
        title={(b) Detalle cerca del origen}]
      \addplot[c1, domain=-0.5:0.5] {exp(-1/(x^2+1e-12))};
      \addplot[c3, densely dashed, domain=-0.5:0.5] {0};
  \end{groupplot}
\end{tikzpicture}
\caption{La función $f(x)=\eul^{-1/x^{2}}$ ($f(0)=0$) es de clase
$\Cinf$ y todas sus derivadas se anulan en el origen, de modo que su serie de
Taylor en $0$ es la función nula (trazo discontinuo). En (b) se aprecia lo
extraordinariamente plana que es $f$ cerca de $0$: en $x=\pm0{,}25$ vale
apenas $\eul^{-16}\approx1{,}1\cdot10^{-7}$.}
\label{fig:suave}
\end{figure}


\begin{observacion}[La explicación está en el plano complejo]\label{obs:suave-complejo}
Sobre el eje imaginario, $z=\iu t$, la misma expresión vale
$\eul^{-1/z^{2}}=\eul^{+1/t^{2}}$, que diverge de forma explosiva cuando
$t\to0$. Vista como función de variable compleja, $\eul^{-1/z^2}$ tiene en
$z=0$ una \emph{singularidad esencial} (\cref{sec:polos}), y ninguna serie
de potencias centrada en $0$ puede representarla. La regularidad sobre la
recta real es engañosa: la analiticidad real es una propiedad que sólo se
entiende cabalmente en $\C$.
\end{observacion}

% -------------------------------------------------------------------------
\subsection{Las soluciones se comportan mejor que los coeficientes}\label{sec:regularidad}
% -------------------------------------------------------------------------
Una idea guía de todo el tema es que \emph{integrar mejora la regularidad y
derivar la empeora}. En los puntos donde $P\neq0$, despejando en
\eqref{eq:edo},
\begin{equation}\label{eq:despeje}
  y''=-\frac{A}{P}\,y'-\frac{B}{P}\,y .
\end{equation}
Para recuperar $y$ hay que integrar dos veces, de modo que $y$ es, en un
sentido heurístico, «dos grados más regular» que los coeficientes $A/P$ y
$B/P$. El ejemplo extremo lo proporciona la delta de Dirac, que ni
siquiera es una función: en el sentido de las distribuciones,
\begin{equation}\label{eq:delta}
  y''=\delta(x)\;\Longrightarrow\;y'=\Heav(x)\ \text{(escalón)}
  \;\Longrightarrow\;y=x\,\Heav(x)\ \text{(rampa, continua)},
\end{equation}
donde se han tomado nulas las constantes de integración. Cada integración
«suaviza»: $\int\delta\,\dif x$ ya es una función acotada (con un salto), y
la siguiente integral es continua (con un pico); véase la \cref{fig:regularidad}.

% Figura: delta -> escalón -> rampa (cada integración suaviza)
\begin{figure}[ht]
\centering
\begin{tikzpicture}[x=0.95cm, y=0.95cm]
  % ---------- Panel 1: delta ----------
  \begin{scope}
    \draw[eje] (-1.7,0) -- (1.9,0) node[right, etiqgris] {$x$};
    \draw[eje] (0,-0.35) -- (0,2.4);
    \draw[-{Stealth[length=7pt,width=6pt]}, line width=1.4pt, cian] (0,0) -- (0,1.75);
    \draw[line width=1.2pt, cian] (-1.6,0) -- (1.75,0);
    \node[etiq, anchor=west] at (0.15,1.6) {$\delta(x)$};
    \node[etiqgris] at (0.05,-0.55) {$0$};
    \node[font=\small, text=gris] at (0.1,-1.25) {singular};
  \end{scope}
  \draw[flecha, line width=0.8pt, ambar] (2.25,1) -- node[above, etiqgris] {$\displaystyle\int$} (3.25,1);
  % ---------- Panel 2: escalón ----------
  \begin{scope}[xshift=5.2cm]
    \draw[eje] (-1.7,0) -- (1.9,0) node[right, etiqgris] {$x$};
    \draw[eje] (0,-0.35) -- (0,2.4);
    \draw[line width=1.4pt, magenta] (-1.6,0) -- (0,0);
    \draw[line width=1.4pt, magenta] (0,1.2) -- (1.75,1.2);
    \draw[guia] (0,0) -- (0,1.2);
    \fill[fondo] (0,0) circle (2.4pt); \draw[magenta, line width=1pt] (0,0) circle (2.4pt);
    \fill[magenta] (0,1.2) circle (2.4pt);
    \node[etiq, anchor=south] at (1.0,1.3) {$\Heav(x)$};
    \node[etiqgris, anchor=east] at (-0.08,1.2) {$1$};
    \node[etiqgris] at (0.05,-0.55) {$0$};
    \node[font=\small, text=gris] at (0.1,-1.25) {acotada, discontinua};
  \end{scope}
  \draw[flecha, line width=0.8pt, ambar] (7.45,1) -- node[above, etiqgris] {$\displaystyle\int$} (8.45,1);
  % ---------- Panel 3: rampa ----------
  \begin{scope}[xshift=10.4cm]
    \draw[eje] (-1.7,0) -- (1.9,0) node[right, etiqgris] {$x$};
    \draw[eje] (0,-0.35) -- (0,2.4);
    \draw[line width=1.4pt, ambar] (-1.6,0) -- (0,0) -- (1.6,1.6);
    \node[etiq, anchor=east] at (-0.2,1.55) {$x\,\Heav(x)$};
    \node[etiqgris] at (0.05,-0.55) {$0$};
    \node[font=\small, text=gris] at (0.1,-1.25) {continua};
  \end{scope}
\end{tikzpicture}
\caption{Cada integración gana un grado de regularidad: la delta de Dirac
(singular) se integra en el escalón de Heaviside $\Heav$ (acotado pero
discontinuo), y éste en la rampa $x\,\Heav(x)$ (continua, con derivada
discontinua). Por la \eqref{eq:despeje}, la solución $y$ se obtiene
integrando dos veces una expresión en los coeficientes.}
\label{fig:regularidad}
\end{figure}


\begin{observacion}[Singularidades fijas]\label{obs:fijas}
Para ecuaciones \emph{lineales}, las únicas singularidades posibles de las
soluciones están en los puntos donde los coeficientes de la forma normal
\eqref{eq:despeje} son singulares, es decir, en los ceros de $P$; por eso se
llaman \emph{singularidades fijas}. Las ecuaciones no lineales no tienen esta
propiedad: la solución de $y'=y^{2}$, $y(0)=1$, es $y=1/(1-x)$, que
explota en $x=1$ aunque el «coeficiente» $y^{2}$ sea perfectamente regular
\cite[cap.~3]{BENDER}. El \cref{teo:fuchs} hará precisa esta idea en el
contexto analítico.
\end{observacion}

% =========================================================================
\section{Singularidades aisladas: polos}\label{sec:polos}
% =========================================================================

Los puntos donde $P$ se anula son aquellos en los que los coeficientes de la
forma normal \eqref{eq:despeje} pueden dejar de ser analíticos. Para
clasificarlos necesitamos medir \emph{cómo} diverge una función. La
herramienta adecuada procede de la variable compleja, y por eso adoptamos
desde ahora el punto de vista de que $x$ puede tomar valores en $\C$.

% -------------------------------------------------------------------------
\subsection{Definición y caracterizaciones}\label{sec:def-polo}
% -------------------------------------------------------------------------

\emph{Idea intuitiva.} Un polo es un punto donde la función se va a
infinito \emph{como una potencia}, igual que $1/x$ o $1/x^{2}$ en $x=0$. El
\emph{orden} mide lo rápido que diverge: $1/x^{2}$ explota más deprisa que
$1/x$.

\begin{definicion}[Polo de orden $p$]\label{def:polo}
Sea $f$ analítica en un entorno reducido
$U^{*}=\{x:0<\abs{x-x_0}<\rho\}$ de $x_0$. Se dice que $f$ tiene un
\emph{polo de orden $p$} en $x_0$, con $p\in\N$, $p\ge1$, si $f$ no admite
extensión analítica a $x_0$ pero
\begin{equation}\label{eq:polo}
  g(x)\defeq(x-x_0)^{p}f(x)\quad\text{es analítica en } x_0
  \quad\text{y}\quad g(x_0)\neq0 .
\end{equation}
Si $f$ coincide en $U^{*}$ con una función analítica en $x_0$, se dice que
$x_0$ es una \emph{singularidad evitable} (o «polo de orden $0$»).
\end{definicion}

Multiplicar por $(x-x_0)^p$ «compensa» exactamente la divergencia.
Despejando, $f=g/(x-x_0)^p$, y como $g(x)\to g(x_0)\eqqcolon c\neq0$,
\begin{equation}\label{eq:asintotico}
  f(x)\sim\frac{c}{(x-x_0)^{p}},\qquad x\to x_0 .
\end{equation}

\begin{proposicion}[Caracterizaciones del orden]\label{prop:orden}
Sea $f$ analítica en un entorno reducido de $x_0$ y no extensible
analíticamente a $x_0$. Para $p\ge1$, son equivalentes:
\begin{enumerate}[label=\textup{(\alph*)}]
  \item $f$ tiene un polo de orden $p$ en $x_0$;
  \item $p$ es el \emph{menor} entero $m\ge0$ tal que $(x-x_0)^{m}f$ es
        analítica en $x_0$;
  \item $f$ admite en un entorno reducido de $x_0$ un desarrollo de
        \emph{Laurent} con un número finito de potencias negativas, la más
        negativa de las cuales es $(x-x_0)^{-p}$:
        \begin{equation}\label{eq:laurent}
          f(x)=\ubr{\frac{g_0}{(x-x_0)^{p}}+\frac{g_1}{(x-x_0)^{p-1}}+\dots
          +\frac{g_{p-1}}{x-x_0}}{\text{parte singular}}
          +\ubr{g_p+g_{p+1}(x-x_0)+\cdots}{\text{parte analítica}}
        \end{equation}
        con $g_0\neq0$.
\end{enumerate}
\end{proposicion}

\begin{proof}
(a)$\Rightarrow$(c). Escribiendo $g=\sum_{j\ge0}g_j(x-x_0)^{j}$ con
$g_0=g(x_0)\neq0$ y dividiendo por $(x-x_0)^{p}$ se obtiene \eqref{eq:laurent}.
(c)$\Rightarrow$(a). Multiplicando \eqref{eq:laurent} por $(x-x_0)^p$ se
obtiene una serie de potencias convergente $g$ con $g(x_0)=g_0\neq0$.

(a)$\Rightarrow$(b). Si $(x-x_0)^{m}f\eqqcolon h$ fuese analítica con
$m<p$, entonces $g=(x-x_0)^{p-m}h$ se anularía en $x_0$, en contra de
$g(x_0)\neq0$. (b)$\Rightarrow$(a). Sea $m$ el mínimo; $m\ge1$ porque $f$
no es extensible. Si $h=(x-x_0)^{m}f$ cumpliera $h(x_0)=0$, se podría
extraer un factor, $h=(x-x_0)h_1$ con $h_1$ analítica, y
$(x-x_0)^{m-1}f=h_1$ sería analítica, contra la minimalidad. Luego
$h(x_0)\neq0$ y $m=p$.
\end{proof}

\begin{observacion}[Por qué «el menor $p$»]\label{obs:menor-p}
Si $(x-x_0)^{2}f$ es analítica, también lo son $(x-x_0)^{3}f$,
$(x-x_0)^{4}f,\dots$, por ser productos de funciones analíticas: multiplicar
de más siempre «funciona». Por eso el orden es el \emph{mínimo} exponente que
lo consigue, y exigir $g(x_0)\neq0$ en \eqref{eq:polo} es una forma
equivalente de imponer esa minimalidad. Un polo de orden $p$ es, en suma,
una serie de potencias a la que se le permiten potencias negativas, pero
sólo hasta $(x-x_0)^{-p}$.
\end{observacion}

% -------------------------------------------------------------------------
\subsection{Ejemplos y clasificación}\label{sec:ejemplos-polos}
% -------------------------------------------------------------------------

\begin{ejemplo}[Polos simples y dobles]\label{ej:polos}
\leavevmode
\begin{enumerate}[label=(\roman*)]
  \item $f(x)=1/x$ en $x_0=0$: $x\,f=1$ es analítica y no nula, luego polo
        de orden $1$ (\emph{polo simple}).
  \item $f(x)=\dfrac{1}{x^{2}(x-1)}$ tiene dos polos de órdenes distintos.
        \begin{itemize}
          \item En $x_0=0$: $x\,f=\dfrac1{x(x-1)}$ sigue divergiendo, pero
                \begin{equation}\label{eq:ej-doble}
                  x^{2}f=\frac{1}{x-1}\ \text{es analítica en }0,\qquad g(0)=-1\neq0,
                \end{equation}
                luego es un polo de orden $2$, con $f\sim-1/x^{2}$ cerca de $0$.
          \item En $x_0=1$: $(x-1)\,f=1/x^{2}$ es analítica en $1$ y vale
                $1\neq0$: polo de orden $1$, con $f\sim1/(x-1)$.
          \item En cualquier otro punto (por ejemplo, $x_0=2$) $f$ es
                analítica: no hay polo.
        \end{itemize}
  \item $f(x)=\dfrac{\sin x}{x^{3}}$ en $x_0=0$. Parece de orden $3$, pero
        \begin{equation}\label{eq:sen-x3}
          \frac{\sin x}{x^{3}}=\frac{x-\frac{x^{3}}{6}+\frac{x^{5}}{120}-\cdots}{x^{3}}
          =\frac1{x^{2}}-\frac16+\frac{x^{2}}{120}-\cdots ,
        \end{equation}
        de modo que $x^{2}f=1-\tfrac{x^{2}}{6}+\cdots$ es analítica con valor
        $1\neq0$, mientras que $x\,f\sim1/x$ todavía diverge: polo de orden
        $2$. El cero del numerador «se come» una potencia.
  \item $f(x)=\dfrac{\sin x}{x}=1-\dfrac{x^{2}}{6}+\cdots$ no tiene un polo:
        definiendo $f(0)=1$ es analítica (singularidad evitable).
\end{enumerate}
\end{ejemplo}

\begin{proposicion}[Regla práctica para cocientes de polinomios]\label{prop:regla}
Sean $N$, $D$ polinomios sin factores comunes salvo, eventualmente, una
potencia de $(x-x_0)$. Si $x_0$ es un cero de multiplicidad $m$ de $D$ y de
multiplicidad $k\ge0$ de $N$, entonces $N/D$ tiene en $x_0$ un polo de orden
$m-k$ si $m>k$, y una singularidad evitable si $m\le k$.
\end{proposicion}

\begin{proof}
Escribiendo $N=(x-x_0)^{k}\tilde N$ y $D=(x-x_0)^{m}\tilde D$ con
$\tilde N(x_0)\tilde D(x_0)\neq0$, se tiene
$(x-x_0)^{m-k}\,N/D=\tilde N/\tilde D$, analítica y no nula en $x_0$
(\cref{prop:algebra}).
\end{proof}

No toda singularidad es un polo. Las funciones $\sqrt{x}$, $\ln x$ y
$\eul^{1/x}$ en $x_0=0$ no lo son, porque ninguna potencia $x^{p}$ las vuelve
analíticas: $x^{p}\sqrt{x}=x^{p+1/2}$ tiene exponente no entero; $x^{p}\ln x$
conserva el logaritmo; y
\[
  \eul^{1/x}=\sum_{m=0}^{\infty}\frac{1}{m!\,x^{m}}
\]
tiene \emph{infinitas} potencias negativas. Son singularidades «peores» que
un polo. La \cref{tab:singularidades} resume la clasificación.

\begin{table}[ht]
\centering
\caption{Clasificación de las singularidades en $x_0$ según el desarrollo de
Laurent $f=\sum_{j}c_j(x-x_0)^{j}$ y el comportamiento de $\abs{f}$.}
\label{tab:singularidades}
\renewcommand{\arraystretch}{1.4}
\small
\begin{tabular}{@{}l l l l@{}}
\toprule
\textcolor{cian}{\bfseries Tipo} & \textcolor{cian}{\bfseries Laurent} &
\textcolor{cian}{\bfseries Comportamiento} & \textcolor{cian}{\bfseries Ejemplo}\\
\midrule
Evitable              & sin potencias negativas            & $f$ acotada               & $\dfrac{\sin x}{x}$\\[4pt]
Polo de orden $p$     & la más negativa es $(x-x_0)^{-p}$  & $\abs{f}\sim\abs{c}\,\abs{x-x_0}^{-p}$ & $\dfrac{1}{x^{2}(x-1)}$\\[4pt]
Esencial              & infinitas potencias negativas      & sin límite en $\C$        & $\eul^{1/x}$\\[4pt]
Punto de ramificación & no existe desarrollo de Laurent    & $f$ multivaluada en $\C$  & $\sqrt{x}$, $\ln x$\\
\bottomrule
\end{tabular}
\end{table}

La \cref{fig:fases} muestra la firma de cada polo en el plano complejo: en
un \emph{retrato de fase} se colorea cada punto $z$ según el argumento
$\arg f(z)$, y alrededor de un polo de orden $p$ el círculo cromático se
recorre exactamente $p$ veces (en sentido opuesto al de un cero de orden $p$),
de acuerdo con el principio del argumento \cite[cap.~4]{AHLFORS}.

% Figura: retratos de fase de 1/z, 1/z^2 y 1/(z^2 (z-1))
% Se usa sombreado plano (una celda = un color) para no interpolar a través
% del salto de +180 a -180 grados del argumento.
\begin{figure}[ht]
\centering
\begin{tikzpicture}
  \begin{groupplot}[group style={group size=3 by 1, horizontal sep=0.55cm},
      oscuro, grid=none, view={0}{90}, width=5.15cm, height=5.15cm,
      axis equal image, colormap name=fase, shader=flat corner,
      point meta min=-180, point meta max=180,
      xmin=-1.5, xmax=1.5, ymin=-1.5, ymax=1.5,
      xtick={-1,0,1}, ytick={-1,0,1},
      samples=\resF, samples y=\resF, domain=-1.5:1.5, y domain=-1.5:1.5,
      tick label style={font=\tiny}, title style={yshift=-2pt},
      xlabel={$\Real z$}, xlabel style={font=\scriptsize, yshift=4pt}]
    % (a) 1/z: arg = -atan2(y,x)
    \nextgroupplot[title={(a) $1/z$\,: polo simple},
        ylabel={$\Imag z$}, ylabel style={font=\scriptsize, yshift=-4pt}]
      \addplot3[surf, point meta={-atan2(y,x)}] (x, y, 0);
      \addplot3[only marks, mark=x, mark size=3pt, line width=1.2pt, blanco]
        coordinates {(0,0,0)};
    % (b) 1/z^2: arg = -2 atan2(y,x), reducido a (-180,180]
    \nextgroupplot[title={(b) $1/z^{2}$\,: polo doble}, yticklabels={}]
      \addplot3[surf, point meta={mod(-2*atan2(y,x)+540,360)-180}] (x, y, 0);
      \addplot3[only marks, mark=x, mark size=3pt, line width=1.2pt, blanco]
        coordinates {(0,0,0)};
    % (c) 1/(z^2 (z-1)): arg = -2 atan2(y,x) - atan2(y,x-1)
    \nextgroupplot[title={(c) $1/\bigl(z^{2}(z-1)\bigr)$}, yticklabels={},
        colorbar right, colorbar style={width=0.22cm, ytick={-180,0,180},
          yticklabels={$-\pi$,$0$,$\pi$}, tick label style={font=\tiny, text=gris},
          axis line style={grismedio}, ylabel={$\arg f$},
          ylabel style={font=\scriptsize, text=texto, yshift=-2pt}}]
      \addplot3[surf, point meta={mod(-2*atan2(y,x)-atan2(y,x-1)+900,360)-180}] (x, y, 0);
      \addplot3[only marks, mark=x, mark size=3pt, line width=1.2pt, blanco]
        coordinates {(0,0,0) (1,0,0)};
  \end{groupplot}
\end{tikzpicture}
\caption{Retratos de fase: cada punto $z$ se colorea según $\arg f(z)$
(barra de la derecha). Alrededor de un polo de orden $p$ (marcado con
$\times$) el círculo cromático se recorre $p$ veces: una vez en (a), dos en
(b). En (c) conviven un polo doble en $z=0$ y uno simple en $z=1$, que son
los del \cref{ej:polos}\,(ii).}
\label{fig:fases}
\end{figure}


% -------------------------------------------------------------------------
\subsection{Relación con la ecuación diferencial}\label{sec:polos-edo}
% -------------------------------------------------------------------------
Escribamos \eqref{eq:edo} en forma normal con
$p\defeq A/P$ y $q\defeq B/P$. Si $P(x_0)\neq0$, $p$ y $q$ son analíticas en
$x_0$ (\cref{prop:algebra}): no hay nada que compensar. Si $P$ tiene un cero
de orden $m$ en $x_0$, por la \cref{prop:regla} los cocientes pueden tener
polos de orden \emph{a lo sumo} $m$ (menos, si $A$ o $B$ se anulan también
en $x_0$ y cancelan factores).

\begin{ejemplo}[Legendre y Bessel]\label{ej:legendre-bessel}
\leavevmode
\begin{enumerate}[label=(\roman*)]
  \item \emph{Legendre}, $P=1-x^{2}$:
        \[
          p=\frac{-2x}{(1-x)(1+x)},\qquad
          q=\frac{\ell(\ell+1)}{(1-x)(1+x)}
        \]
        tienen polos de orden $1$ en $x=\pm1$ (si $\ell(\ell+1)\neq0$).
  \item \emph{Bessel}, $x^{2}y''+x\,y'+(x^{2}-\nu^{2})\,y=0$:
        \[
          p=\frac1x\ \ (\text{orden }1),\qquad
          q=\frac{x^{2}-\nu^{2}}{x^{2}}\ \ (\text{orden }2\text{ si }\nu\neq0)
        \]
        en $x=0$. Si $\nu=0$, $q\equiv1$ es analítica.
\end{enumerate}
\end{ejemplo}

\begin{observacion}[Avance de la Parte II]\label{obs:avance}
El orden de estos polos es lo que distinguirá, en la segunda parte del tema,
los puntos \emph{singulares regulares} ---en los que $p$ tiene a lo sumo un
polo de orden $1$ y $q$ a lo sumo uno de orden $2$, es decir, $(x-x_0)\,p$ y
$(x-x_0)^{2}q$ son analíticas--- de los \emph{irregulares}. En los primeros
se buscan soluciones de la forma $(x-x_0)^{r}\sum a_n(x-x_0)^n$ (método de
Frobenius), con un exponente $r$ fijado por la \emph{ecuación indicial}
\cite{TEMAV}. Tanto $x=\pm1$ en Legendre como $x=0$ en Bessel son singulares
regulares.
\end{observacion}

% -------------------------------------------------------------------------
\subsection{Por qué los polos complejos limitan el radio}\label{sec:radio}
% -------------------------------------------------------------------------
La función $f(x)=1/(1+x^{2})$ no tiene ninguna singularidad real: es
analítica en todo $\R$, acotada y de clase $\Cinf$. Y, sin embargo, su
desarrollo en $x_0=0$,
\begin{equation}\label{eq:lorentz}
  \frac{1}{1+x^{2}}=1-x^{2}+x^{4}-x^{6}+\cdots=\sum_{n=0}^{\infty}(-1)^{n}x^{2n},
\end{equation}
\emph{sólo converge para $\abs{x}<1$}: es una serie geométrica de razón
$-x^{2}$. La razón está en el plano complejo: $1+z^{2}=0$ en $z=\pm\iu$, que
son polos de orden $1$ situados a distancia $1$ del origen. Una serie de
potencias converge en un \emph{disco}, y ese disco no puede contener un
polo; por eso el radio es la distancia al polo complejo más cercano
(\cref{fig:radio}).

\begin{teorema}[Radio de convergencia y singularidades]\label{teo:radio}
Sea $f$ holomorfa en el disco $\disco{z_0}{\rho}=\{z\in\C:\abs{z-z_0}<\rho\}$.
Entonces la serie de Taylor de $f$ en $z_0$ converge a $f$ en todo
$\disco{z_0}{\rho}$. En particular, si $f$ es racional, el radio de
convergencia de su serie de Taylor en $z_0$ es exactamente la distancia de
$z_0$ al polo más próximo.
\end{teorema}

\begin{proof}
La primera afirmación es el teorema de Cauchy--Taylor \cite[§4.3]{AHLFORS}.
Para la segunda, sea $d$ la distancia al polo más cercano $z_1$; la primera
parte da $R\ge d$. Si fuera $R>d$, la suma de la serie sería continua en el
disco cerrado $\overline{\disco{z_0}{d}}$, y por tanto acotada en él, y
coincidiría con $f$ en $\disco{z_0}{d}$; pero $\abs{f(z)}\to\infty$ cuando
$z\to z_1$ dentro de ese disco. Luego $R=d$.
\end{proof}

% Figura: 1/(1+x^2) — polos complejos ±i y radio de convergencia R = 1
\begin{figure}[ht]
\centering
\begin{tikzpicture}
  % ---------------- (a) Plano complejo: retrato de fase y disco ----------------
  \begin{axis}[name=fase, oscuro, grid=none, view={0}{90},
      width=6.3cm, height=6.3cm, axis equal image,
      colormap name=fase, shader=flat corner,
      point meta min=-180, point meta max=180,
      xmin=-2, xmax=2, ymin=-2, ymax=2,
      xtick={-2,-1,0,1,2}, ytick={-2,-1,0,1,2},
      samples=\resF, samples y=\resF, domain=-2:2, y domain=-2:2,
      xlabel={$\Real z$}, ylabel={$\Imag z$},
      title={(a) Plano complejo: $\arg\dfrac{1}{1+z^{2}}$}]
    \addplot3[surf, point meta={-atan2(2*x*y,1+x^2-y^2)}] (x, y, 0);
    % disco de convergencia |z| < 1
    \addplot3[blanco, line width=1.1pt, domain=0:360, samples=120, samples y=0]
      ({cos(x)}, {sin(x)}, 0);
    % intervalo real de convergencia
    \addplot3[blanco, line width=2.6pt, samples y=0] coordinates {(-1,0,0) (1,0,0)};
    \addplot3[only marks, mark=*, mark size=1.8pt, blanco, samples y=0]
      coordinates {(-1,0,0) (1,0,0)};
    % radio
    \addplot3[blanco, line width=0.8pt, densely dashed, samples y=0]
      coordinates {(0,0,0) (0.7071,-0.7071,0)};
    % polos
    \addplot3[only marks, mark=x, mark size=4pt, line width=1.5pt, blanco, samples y=0]
      coordinates {(0,1,0) (0,-1,0)};
    \node[fill=fondo, fill opacity=0.8, text opacity=1, font=\scriptsize,
          inner sep=1.5pt, rounded corners=1pt] at (axis cs:0.42,1.36,0) {$+\iu$};
    \node[fill=fondo, fill opacity=0.8, text opacity=1, font=\scriptsize,
          inner sep=1.5pt, rounded corners=1pt] at (axis cs:0.42,-1.36,0) {$-\iu$};
    \node[fill=fondo, fill opacity=0.8, text opacity=1, font=\scriptsize,
          inner sep=1.5pt, rounded corners=1pt] at (axis cs:1.35,-0.95,0) {$R=1$};
  \end{axis}

  % ---------------- (b) Recta real: sumas parciales ----------------------------
  \begin{axis}[at={($(fase.east)+(1.45cm,0)$)}, anchor=west, oscuro,
      width=7.2cm, height=5.9cm,
      xmin=-1.6, xmax=1.6, ymin=-0.6, ymax=1.6,
      xtick={-1.5,-1,-0.5,0,0.5,1,1.5}, ytick={-0.5,0,0.5,1,1.5},
      xlabel={$x$}, title={(b) Recta real: sumas parciales $S_N$},
      legend style={at={(0.5,0.03)}, anchor=south, legend columns=2},
      samples=241, domain=-1.6:1.6]
    \fill[grisoscuro, opacity=0.55] (axis cs:-1.6,-0.6) rectangle (axis cs:-1,1.6);
    \fill[grisoscuro, opacity=0.55] (axis cs:1,-0.6) rectangle (axis cs:1.6,1.6);
    \addplot[c2, restrict y to domain=-5:5] {(1+x^6)/(1+x^2)};
    \addplot[c3, restrict y to domain=-5:5] {(1-x^12)/(1+x^2)};
    \addplot[c4, restrict y to domain=-5:5] {(1+x^26)/(1+x^2)};
    \addplot[blanco, line width=1.6pt] {1/(1+x^2)};
    \legend{$S_{4}$, $S_{10}$, $S_{24}$, $\frac{1}{1+x^{2}}$}
  \end{axis}
\end{tikzpicture}
\caption{(a) Retrato de fase de $f(z)=1/(1+z^{2})$. Los polos $z=\pm\iu$
(marcados con $\times$) están a distancia $1$ del origen; el mayor disco
centrado en $0$ que no los contiene es $\abs{z}<1$, cuya intersección con la
recta real es el intervalo $(-1,1)$ (trazo grueso). (b) Las sumas parciales
$S_{2N}(x)=\sum_{n=0}^{N}(-1)^n x^{2n}$ de \eqref{eq:lorentz} se acercan a
$f$ en $\abs{x}<1$ y se separan bruscamente de ella en $\abs{x}>1$
(zonas sombreadas), pese a que $f$ es perfectamente regular sobre $\R$.}
\label{fig:radio}
\end{figure}


Esta es la idea que explota el \cref{teo:fuchs}: para conocer dónde
converge la serie de una solución \emph{real} hay que mirar dónde se anula
$P$ en el plano \emph{complejo}.

% =========================================================================
\section{Puntos ordinarios: soluciones en serie entera}\label{sec:ordinarios}
% =========================================================================

% -------------------------------------------------------------------------
\subsection{Puntos ordinarios y puntos singulares}\label{sec:def-ordinario}
% -------------------------------------------------------------------------

\begin{definicion}[Punto ordinario y punto singular]\label{def:ordinario}
Un punto $x_0\in\R$ es un \emph{punto ordinario} de la ecuación
\eqref{eq:edo} si \emph{no anula} el polinomio $P$, es decir,
$P(x_0)\neq0$. En caso contrario, $P(x_0)=0$, se dice que $x_0$ es un
\emph{punto singular}.
\end{definicion}

\begin{ejemplo}\label{ej:ordinario}
En $(1-x^{2})\,y''-2x\,y'+\tfrac34\,y=0$ el punto $x_0=0$ es ordinario. De
hecho, todo $x_0\neq\pm1$ es ordinario, mientras que $x_0=\pm1$ son
singulares \cite{TEMAV}.
\end{ejemplo}

Si $x_0$ es ordinario podemos dividir por $P$ en un entorno de $x_0$.
Definimos
\begin{equation}\label{eq:pq}
  p(x)\defeq\frac{A(x)}{P(x)},\qquad
  q(x)\defeq\frac{B(x)}{P(x)},\qquad
  z\defeq x-x_0 ,
\end{equation}
con lo que la ecuación adopta la \emph{forma normal}
\begin{equation}\label{eq:normal}
  y''+p\,y'+q\,y=0 .
\end{equation}
Como $P(x_0)\neq0$, $p$ y $q$ son analíticas en $x_0$ por ser cocientes de
funciones analíticas con denominador no nulo (\cref{prop:algebra}):
\begin{equation}\label{eq:pq-serie}
  p=\sum_{k=0}^{\infty}p_k\,z^{k},\qquad
  q=\sum_{k=0}^{\infty}q_k\,z^{k},\qquad \abs{z}<R\defeq\min(R_p,R_q).
\end{equation}

\begin{observacion}[Una definición más general]\label{obs:ordinario-general}
Lo que interviene en toda la teoría es la analiticidad de $p$ y $q$ en
$x_0$, no el valor de $P(x_0)$. Por eso muchos textos llaman ordinario a
todo punto en el que $p$ y $q$ son analíticas \cite[§5.2]{BOYCE}. Por
ejemplo, en $x\,y''+(\sin x)\,y=0$ se tiene $P(0)=0$, pero
$q=\sin x/x$ tiene en $0$ una singularidad evitable
(\cref{ej:polos}\,(iv)): el punto $x_0=0$ es ordinario en sentido amplio y
todos los resultados de esta sección se le aplican.
\end{observacion}

% -------------------------------------------------------------------------
\subsection{El teorema de Fuchs}\label{sec:fuchs}
% -------------------------------------------------------------------------
El resultado central del tema afirma que en un punto ordinario el método de
series no sólo es formalmente viable, sino que produce soluciones
convergentes, y da una cota inferior explícita del radio
\cite[cap.~16]{RILEY}, \cite[cap.~5]{SIMMONS}.

\begin{teorema}[Fuchs, puntos ordinarios]\label{teo:fuchs}
Sea $x_0$ un punto ordinario de \eqref{eq:edo} y sean $p$, $q$ analíticas
en $\abs{z}<R$, con $z=x-x_0$ y $R=\min(R_p,R_q)$. Entonces:
\begin{enumerate}[label=\textup{(\roman*)}]
  \item para cada par $(a_0,a_1)\in\R^{2}$ existe una única serie
        $y=\sum_{n\ge0}a_n z^{n}$ que es solución de \eqref{eq:normal} con
        $y(x_0)=a_0$, $y'(x_0)=a_1$, y converge \emph{al menos} en
        $\abs{z}<R$;
  \item toda solución de \eqref{eq:normal} en $(x_0-R,\,x_0+R)$ es de esta
        forma; en particular, toda solución es analítica en $x_0$;
  \item si $A$ y $B$ son polinomios (o, más en general, funciones enteras),
        \begin{equation}\label{eq:cota-radio}
          R\;\ge\;\dist\bigl(x_0,\,\mathcal Z(P)\bigr),
          \qquad \mathcal Z(P)\defeq\{\zeta\in\C: P(\zeta)=0\},
        \end{equation}
        es decir, el radio es al menos la distancia, \emph{en el plano
        complejo}, de $x_0$ al cero de $P$ más cercano.
\end{enumerate}
\end{teorema}

La demostración se da en la \cref{sec:demo-fuchs}, una vez obtenida la
relación de recurrencia. La \cref{fig:fuchs} ilustra la cota
\eqref{eq:cota-radio}: es la versión, para las soluciones de la EDO, de lo
que el \cref{teo:radio} afirma para funciones racionales.

% Figura: disco de convergencia garantizado por el teorema de Fuchs
\begin{figure}[ht]
\centering
\begin{tikzpicture}[scale=1.15]
  % coordenadas: x0 = 0.6; ceros de P: zeta1 = 3.3 (real), zeta2,3 = -0.6 ± 1.9 i
  \def\xo{0.6}
  \pgfmathsetmacro{\radio}{sqrt((0.6+0.6)^2+1.9^2)}   % = dist(x0, zeta2) ≈ 2.247
  % disco
  \fill[cian, opacity=0.13] (\xo,0) circle (\radio);
  \draw[cian, line width=0.9pt] (\xo,0) circle (\radio);
  % ejes
  \draw[eje] (-2.6,0) -- (4.5,0) node[right, etiqgris] {$\Real z$};
  \draw[eje] (0,-2.75) -- (0,2.85) node[above, etiqgris] {$\Imag z$};
  % intervalo real de convergencia
  \draw[ambar, line width=2.4pt] ({\xo-\radio},0) -- ({\xo+\radio},0);
  \fill[ambar] ({\xo-\radio},0) circle (1.7pt);
  \fill[ambar] ({\xo+\radio},0) circle (1.7pt);
  \node[font=\footnotesize, text=ambar, anchor=north] at ({\xo-\radio-0.05},-0.12) {$x_0-R$};
  \node[font=\footnotesize, text=ambar, anchor=north] at ({\xo+\radio+0.05},-0.12) {$x_0+R$};
  % radio hasta el cero complejo más cercano
  \draw[blanco, line width=0.8pt, densely dashed] (\xo,0) -- (-0.6,1.9)
     node[midway, right=3pt, font=\small] {$R$};
  % centro
  \node[punto] at (\xo,0) {};
  \node[font=\small, anchor=north] at (\xo+0.12,-0.12) {$x_0$};
  % ceros de P
  \node[polo] at (3.3,0) {};
  \node[polo] at (-0.6,1.9) {};
  \node[polo] at (-0.6,-1.9) {};
  \node[font=\small, text=rojo, anchor=south] at (3.3,0.14) {$\zeta_1$};
  \node[font=\small, text=rojo, anchor=east] at (-0.72,1.98) {$\zeta_2$};
  \node[font=\small, text=rojo, anchor=east] at (-0.72,-1.98) {$\zeta_3=\overline{\zeta_2}$};
  % rótulo del disco
  \node[font=\small, text=cian, anchor=west] at (2.15,1.75) {$\abs{z-x_0}<R$};
  % leyenda
  \begin{scope}[shift={(5.0,1.6)}]
    \node[polo] at (0,0) {};
    \node[etiqgris, anchor=west] at (0.2,0) {ceros de $P$ en $\C$};
    \draw[ambar, line width=2.4pt] (-0.15,-0.55) -- (0.15,-0.55);
    \node[etiqgris, anchor=west] at (0.2,-0.55) {convergencia en $\R$};
  \end{scope}
\end{tikzpicture}
\caption{Teorema de Fuchs. El disco de convergencia garantizado está
centrado en el punto ordinario $x_0$ y llega hasta el cero complejo de $P$
más próximo, aquí el par conjugado $\zeta_{2,3}$. Aunque el único cero
\emph{real}, $\zeta_1$, está más lejos, es $\zeta_2$ quien fija
$R=\abs{\zeta_2-x_0}$ y, con él, el intervalo real $(x_0-R,\,x_0+R)$ en
el que converge la serie.}
\label{fig:fuchs}
\end{figure}


\begin{observacion}[El radio garantizado es una cota inferior]\label{obs:cota}
La cota \eqref{eq:cota-radio} puede no alcanzarse: si una de las series
solución se trunca en un polinomio, ésta converge en todo $\R$ aunque $P$
tenga ceros (polinomios de Legendre, \cref{sec:legendre}). Lo que el teorema
asegura es que \emph{ninguna} solución deja de converger antes de llegar al
cero complejo de $P$ más próximo.
\end{observacion}

% -------------------------------------------------------------------------
\subsection{Ansatz y derivación término a término}\label{sec:ansatz}
% -------------------------------------------------------------------------
Proponemos
\begin{equation}\label{eq:ansatz}
  y=\sum_{n=0}^{\infty}a_n\,z^{n}.
\end{equation}
Derivando término a término (\cref{teo:series}),
\begin{align}
  y'  &=\sum_{n=1}^{\infty}n\,a_n\,z^{n-1}, \label{eq:yp}\\
  y'' &=\sum_{n=2}^{\infty}n(n-1)\,a_n\,z^{n-2}. \label{eq:ypp}
\end{align}
Los límites inferiores suben porque el término $n=0$ de $y'$ (y los
$n=0,1$ de $y''$) son nulos. De momento el cálculo es \emph{formal}: la
convergencia de \eqref{eq:ansatz} se establecerá en la \cref{sec:demo-fuchs},
y es entonces cuando el \cref{teo:series} legitima a posteriori cada paso.

% -------------------------------------------------------------------------
\subsection{Sustitución y cambio de índices}\label{sec:indices}
% -------------------------------------------------------------------------
Introduciendo \eqref{eq:pq-serie}, \eqref{eq:yp} y \eqref{eq:ypp} en
\eqref{eq:normal}:
\begin{equation}\label{eq:sustitucion}
  \sum_{n=2}^{\infty}n(n-1)\,a_n z^{n-2}
  +\Bigl(\sum_{k=0}^{\infty}p_k z^{k}\Bigr)\Bigl(\sum_{n=1}^{\infty}n\,a_n z^{n-1}\Bigr)
  +\Bigl(\sum_{k=0}^{\infty}q_k z^{k}\Bigr)\Bigl(\sum_{n=0}^{\infty}a_n z^{n}\Bigr)=0 .
\end{equation}
Queremos escribir todas las series en potencias $z^{n}$. El índice de
suma es \emph{mudo} y puede desplazarse libremente:
\begin{equation}\label{eq:desplazamiento}
  \sum_{n=m}^{\infty}c_n\,z^{n-m}
  \;\overset{n\,\mapsto\,n+m}{=}\;
  \sum_{n=0}^{\infty}c_{n+m}\,z^{n}.
\end{equation}
\begin{itemize}
  \item En $y''$ tomamos $m=2$ (es decir, llamamos $n$ a $n-2$):
        \begin{equation}\label{eq:ypp-desplazada}
          y''=\sum_{n=0}^{\infty}(n+2)(n+1)\,a_{n+2}\,z^{n}.
        \end{equation}
  \item En $y'$ tomamos $m=1$:
        \begin{equation}\label{eq:yp-desplazada}
          y'=\sum_{n=0}^{\infty}(n+1)\,a_{n+1}\,z^{n}.
        \end{equation}
\end{itemize}
Ahora las tres series empiezan en $n=0$ y llevan $z^{n}$; usando un índice
$n'$ en los factores para no confundirlo con el de los productos,
\begin{equation}\label{eq:sustitucion2}
  \sum_{n}(n+2)(n+1)a_{n+2}z^{n}
  +\Bigl(\sum_{k}p_k z^{k}\Bigr)\Bigl(\sum_{n'}(n'+1)a_{n'+1}z^{n'}\Bigr)
  +\Bigl(\sum_{k}q_k z^{k}\Bigr)\Bigl(\sum_{n'}a_{n'}z^{n'}\Bigr)=0 .
\end{equation}

% -------------------------------------------------------------------------
\subsection{Producto de Cauchy}\label{sec:cauchy}
% -------------------------------------------------------------------------

\begin{lema}[Producto de Cauchy]\label{lem:cauchy}
Si las series $\sum_{k\ge0}c_kz^{k}$ y $\sum_{n'\ge0}b_{n'}z^{n'}$ convergen
absolutamente en $\abs{z}<R$, entonces, en ese disco,
\begin{equation}\label{eq:cauchy}
  \Bigl(\sum_{k=0}^{\infty}c_k z^{k}\Bigr)\Bigl(\sum_{n'=0}^{\infty}b_{n'}z^{n'}\Bigr)
  =\sum_{n=0}^{\infty}\Bigl(\sum_{k=0}^{n}c_k\,b_{n-k}\Bigr)z^{n},
\end{equation}
y la serie del miembro derecho converge también absolutamente.
\end{lema}

\begin{proof}
Al multiplicar aparecen los términos $c_kb_{n'}z^{k}z^{n'}=c_kb_{n'}z^{k+n'}$
para todos los pares $(k,n')\in\N^{2}$. La serie doble correspondiente
converge absolutamente, pues
$\sum_{k,n'}\abs{c_k}\abs{b_{n'}}\abs{z}^{k+n'}
=\bigl(\sum_k\abs{c_k}\abs{z}^k\bigr)\bigl(\sum_{n'}\abs{b_{n'}}\abs{z}^{n'}\bigr)<\infty$,
y por tanto puede sumarse en cualquier orden. Agrupamos los términos con la
misma potencia total $n=k+n'$, es decir, recorremos las diagonales de la
\cref{fig:cauchy}. Fijado $n$, se tiene $n'=n-k$, y como $n'\ge0$, el
índice $k$ sólo recorre $0,1,\dots,n$.
\end{proof}

% Figura: producto de Cauchy como suma por diagonales k + n' = n
\begin{figure}[ht]
\centering
\begin{tikzpicture}[x=0.82cm, y=0.82cm]
  % diagonales k + n' = n
  \foreach \n in {1,...,6}{
    \draw[guia] (\n+0.3,-0.3) -- (-0.3,\n+0.3);}
  % diagonal destacada n = 4
  \draw[cian, line width=1.2pt] (4.35,-0.35) -- (-0.35,4.35);
  % ejes
  \draw[eje] (-0.4,0) -- (6.6,0) node[right, etiqgris] {$k$};
  \draw[eje] (0,-0.4) -- (0,6.6) node[above, etiqgris] {$n'$};
  \foreach \t in {0,...,5}{
    \node[etiqgris, anchor=north] at (\t,-0.45) {$\t$};
    \node[etiqgris, anchor=east] at (-0.45,\t) {$\t$};}
  % malla de pares (k, n')
  \foreach \k in {0,...,5}{\foreach \m in {0,...,5}{
    \fill[grismedio] (\k,\m) circle (2.2pt);}}
  \foreach \k in {0,...,4}{
    \pgfmathtruncatemacro{\m}{4-\k}
    \fill[cian] (\k,\m) circle (3.4pt);}
  % rótulo de la diagonal
  \node[nodocian, anchor=west, align=left] (r) at (7.4,4.1)
    {\textcolor{cian}{\bfseries Diagonal $n=4$}\\[2pt]
     $c_0b_4+c_1b_3+c_2b_2+c_3b_1+c_4b_0$\\[1pt]
     \footnotesize\color{gris} coeficiente de $z^{4}$ en el producto};
  \draw[flecha, cian] (r.west) to[out=180, in=45] (2.15,2.15);
  \node[nodo, anchor=west, align=left] at (7.4,1.4)
    {\footnotesize\color{gris} Cada punto $(k,n')$ es el término\\
     \footnotesize\color{gris} $c_k b_{n'}z^{k+n'}$; en la diagonal\\
     \footnotesize\color{gris} $k+n'=n$ todos llevan $z^{n}$.};
\end{tikzpicture}
\caption{Producto de Cauchy de dos series de potencias. Los términos
$c_kb_{n'}z^{k+n'}$ se agrupan por diagonales $k+n'=n$; como $n'=n-k\ge0$,
en la diagonal $n$ el índice $k$ sólo recorre $0,1,\dots,n$, lo que da el
coeficiente $\sum_{k=0}^{n}c_kb_{n-k}$ de \eqref{eq:cauchy}.}
\label{fig:cauchy}
\end{figure}


Aplicando el \cref{lem:cauchy} con $b_{n'}=(n'+1)\,a_{n'+1}$, de modo que
$b_{n-k}=(n-k+1)\,a_{n-k+1}$, y con $b_{n'}=a_{n'}$:
\begin{align}
  p\,y' &=\sum_{n=0}^{\infty}\Bigl(\sum_{k=0}^{n}(n-k+1)\,p_k\,a_{n-k+1}\Bigr)z^{n},\label{eq:pyp}\\
  q\,y  &=\sum_{n=0}^{\infty}\Bigl(\sum_{k=0}^{n}q_k\,a_{n-k}\Bigr)z^{n}.\label{eq:qy}
\end{align}

% -------------------------------------------------------------------------
\subsection{Identificación de coeficientes: la relación de recurrencia}\label{sec:recurrencia}
% -------------------------------------------------------------------------
Toda la ecuación es ahora una única serie $\sum_n C_n z^{n}=0$ válida para
todo $\abs{z}<R$. Por la unicidad del desarrollo en serie de potencias
(\cref{cor:unicidad}) ---la serie nula sólo tiene coeficientes nulos---,
cada coeficiente debe anularse.

\begin{resultado}[Relación de recurrencia en un punto ordinario]\label{res:recurrencia}
Los coeficientes de toda solución \eqref{eq:ansatz} de \eqref{eq:normal}
satisfacen, para todo $n\ge0$,
\begin{equation}\label{eq:recurrencia}
  (n+2)(n+1)\,a_{n+2}
  +\sum_{k=0}^{n}\Bigl[(n-k+1)\,p_k\,a_{n-k+1}+q_k\,a_{n-k}\Bigr]=0 .
\end{equation}
\end{resultado}

Ésta es la ecuación en diferencias anunciada en la \cref{sec:por-que}: la
EDO se ha convertido en una relación de recurrencia para los coeficientes.
Si $p$ y $q$ son polinomios de grado $\le d$, la suma en \eqref{eq:recurrencia}
tiene a lo sumo $d+1$ términos no nulos y la recurrencia es de orden finito;
en general involucra todos los coeficientes anteriores.

\begin{observacion}[Erratas frecuentes]\label{obs:erratas}
Al transcribir \eqref{eq:recurrencia} conviene vigilar dos detalles en los
que es fácil equivocarse: en el término de $q$ la potencia que acompaña a
$q_k$ es $z^{k}$, no $z^{2}$; y en el de $p$ el subíndice correcto es
$a_{n-k+1}$, multiplicado por el factor $(n-k+1)$ procedente de la derivada.
\end{observacion}

% -------------------------------------------------------------------------
\subsection{Resolución de la recurrencia}\label{sec:resolucion}
% -------------------------------------------------------------------------
Como $(n+2)(n+1)\neq0$ para $n\ge0$, podemos despejar el coeficiente de
mayor índice:
\begin{equation}\label{eq:despeje-an2}
  \boxed{\;a_{n+2}=-\frac{1}{(n+2)(n+1)}
  \sum_{k=0}^{n}\Bigl[(n-k+1)\,p_k\,a_{n-k+1}+q_k\,a_{n-k}\Bigr],
  \qquad n\ge0.\;}
\end{equation}
El miembro derecho sólo contiene $a_0,a_1,\dots,a_{n+1}$: conocidos todos los
anteriores, $a_{n+2}$ queda determinado (\cref{fig:dependencia}).

\begin{proposicion}[Determinación por inducción]\label{prop:induccion}
Para cada par $(a_0,a_1)$ existe una única sucesión $(a_n)_{n\ge0}$ que
cumple \eqref{eq:despeje-an2}. Además, cada $a_n$ es una combinación lineal
de $a_0$ y $a_1$.
\end{proposicion}

\begin{proof}
Por inducción fuerte sobre $n$: \eqref{eq:despeje-an2} define $a_{n+2}$ a
partir de $a_0,\dots,a_{n+1}$, ya determinados. La linealidad se debe a que
\eqref{eq:despeje-an2} es lineal en los $a_j$.
\end{proof}

% Figura: estructura de dependencia de la recurrencia general
\begin{figure}[ht]
\centering
\begin{tikzpicture}
  \tikzset{coef/.style={circle, draw=grismedio, fill=cajafondo, line width=0.6pt,
                        minimum size=0.95cm, inner sep=0pt, font=\small}}
  \foreach \i in {0,...,6}{
    \node[coef] (a\i) at (1.75*\i,0) {$a_{\i}$};}
  \node[font=\small, text=gris] at (1.75*7-0.35,0) {$\cdots$};
  % datos libres
  \node[coef, draw=ambar, line width=1pt] at (a0) {$a_{0}$};
  \node[coef, draw=ambar, line width=1pt] at (a1) {$a_{1}$};
  % coeficiente calculado
  \node[coef, draw=cian, line width=1.2pt] at (a4) {$a_{4}$};
  % flechas de dependencia hacia a4 (n = 2)
  \foreach \j/\h in {0/62, 1/52, 2/40, 3/30}{
    \draw[flecha, cian!75] (a\j.north) to[out=\h, in=180-\h+8] (a4.north);}
  % llave bajo los datos libres
  \draw[decorate, decoration={brace, amplitude=5pt, mirror}, ambar, line width=0.6pt]
    ($(a0.south west)+(0,-0.15)$) -- ($(a1.south east)+(0,-0.15)$)
    node[midway, below=7pt, font=\footnotesize, text=ambar, align=center]
    {datos libres\\$a_0=y(x_0)$,\ $a_1=y'(x_0)$};
  \draw[decorate, decoration={brace, amplitude=5pt, mirror}, grismedio, line width=0.6pt]
    ($(a2.south west)+(0,-0.15)$) -- ($(a6.south east)+(0,-0.15)$)
    node[midway, below=7pt, font=\footnotesize, text=gris, align=center]
    {determinados por \eqref{eq:despeje-an2},\\uno tras otro};
  \node[font=\footnotesize, text=cian, anchor=west, align=left] at ($(a5.north)+(-0.1,0.95)$)
    {$a_{4}$ depende de $a_0,a_1,a_2,a_3$\\($n=2$ en la recurrencia)};
\end{tikzpicture}
\caption{Estructura de la recurrencia \eqref{eq:despeje-an2}: el coeficiente
$a_{n+2}$ depende de \emph{todos} los anteriores (aquí, $a_4$). Sólo $a_0$ y
$a_1$ quedan libres; son las dos constantes de integración.}
\label{fig:dependencia}
\end{figure}


\paragraph{Primeros términos}
Particularizando \eqref{eq:despeje-an2}:
\begin{itemize}
  \item $n=0$ (sólo $k=0$):
        \begin{equation}\label{eq:a2}
          a_2=-\tfrac12\bigl(p_0a_1+q_0a_0\bigr).
        \end{equation}
  \item $n=1$ ($k=0,1$):
        \begin{equation}\label{eq:a3}
          a_3=-\tfrac16\bigl(2p_0a_2+q_0a_1+p_1a_1+q_1a_0\bigr).
        \end{equation}
  \item $n=2$ ($k=0,1,2$):
        \begin{equation}\label{eq:a4}
          a_4=-\tfrac1{12}\bigl(3p_0a_3+(q_0+2p_1)\,a_2+(q_1+p_2)\,a_1+q_2a_0\bigr).
        \end{equation}
\end{itemize}

\begin{observacion}[Comprobación]\label{obs:comprobacion}
La fórmula \eqref{eq:a2} puede verificarse sin series: evaluando
\eqref{eq:normal} en $x_0$ se obtiene
$y''(x_0)=-p(x_0)\,y'(x_0)-q(x_0)\,y(x_0)=-p_0a_1-q_0a_0$, y por el
\cref{teo:series}\,(iii), $a_2=y''(x_0)/2$. Del mismo modo, derivando la
ecuación sucesivamente se recupera cada $a_n=y^{(n)}(x_0)/n!$: la
recurrencia es el algoritmo de Taylor organizado.
\end{observacion}

% -------------------------------------------------------------------------
\subsection{Condiciones iniciales y sistema fundamental}\label{sec:fundamental}
% -------------------------------------------------------------------------
Evaluando la serie \eqref{eq:ansatz} y su derivada \eqref{eq:yp} en $z=0$,
\begin{equation}\label{eq:a0a1}
  a_0=y(x_0),\qquad a_1=y'(x_0).
\end{equation}
Son las condiciones iniciales: las dos constantes de integración que se
esperan en una EDO de segundo orden. La recurrencia no las fija; quedan
libres.

Como la ecuación es lineal, basta elegir dos juegos de condiciones
iniciales:
\begin{equation}\label{eq:y1y2}
  y_1\defeq y\big|_{(a_0,a_1)=(1,0)},\qquad
  y_2\defeq y\big|_{(a_0,a_1)=(0,1)}.
\end{equation}

\begin{proposicion}[Sistema fundamental]\label{prop:fundamental}
Las funciones $y_1$, $y_2$ forman un sistema fundamental de soluciones de
\eqref{eq:normal} en $(x_0-R,\,x_0+R)$, y la solución general es
\begin{equation}\label{eq:general}
  y=a_0\,y_1+a_1\,y_2,\qquad a_0,a_1\in\R .
\end{equation}
\end{proposicion}

\begin{proof}
Su wronskiano en $x_0$ es
\begin{equation}\label{eq:wronskiano}
  W(x_0)=\begin{vmatrix} y_1(x_0) & y_2(x_0)\\ y_1'(x_0) & y_2'(x_0)\end{vmatrix}
  =\begin{vmatrix}1&0\\0&1\end{vmatrix}=1\neq0 .
\end{equation}
Además $W'=y_1y_2''-y_1''y_2=-p\,W$, de donde (fórmula de Abel)
$W(x)=\exp\bigl(-\int_{x_0}^{x}p(t)\,\dif t\bigr)\neq0$ en todo el
intervalo: $y_1$ e $y_2$ son linealmente independientes. Por la
\cref{prop:induccion}, la solución con datos $(a_0,a_1)$ tiene coeficientes
lineales en $(a_0,a_1)$, luego es $a_0y_1+a_1y_2$.
\end{proof}

% -------------------------------------------------------------------------
\subsection{Demostración del teorema de Fuchs}\label{sec:demo-fuchs}
% -------------------------------------------------------------------------
La única afirmación no trivial del \cref{teo:fuchs} es que la serie formal
construida con \eqref{eq:despeje-an2} \emph{converge}. Lo probaremos con el
\emph{método de las series mayorantes} de Cauchy: acotaremos $\abs{a_n}$ por
una sucesión $b_n$ cuyo comportamiento conocemos exactamente
\cite[cap.~3]{CODDINGTON}, \cite[§4.1]{TESCHL}.

\begin{proof}[Demostración del \cref{teo:fuchs}]
Si $a_0=a_1=0$, todos los $a_n$ son nulos y no hay nada que probar;
supongamos $(a_0,a_1)\neq(0,0)$.

\emph{Paso 1: cotas de Cauchy.} Fijemos $0<r<R$. Como $\sum p_kr^{k}$ y
$\sum q_kr^{k}$ convergen, sus términos están acotados: existe $M>0$ tal que
\begin{equation}\label{eq:cotas-cauchy}
  \abs{p_k}\le M\,r^{-k},\qquad \abs{q_k}\le M\,r^{-k}\qquad\forall k\ge0 .
\end{equation}

\emph{Paso 2: sucesión mayorante.} Definimos $b_0=\abs{a_0}$,
$b_1=\abs{a_1}$ y, para $n\ge0$,
\begin{equation}\label{eq:mayorante}
  (n+2)(n+1)\,b_{n+2}=M\sum_{j=0}^{n}r^{\,j-n}\bigl[(j+1)\,b_{j+1}+b_j\bigr]+M\,r\,b_{n+1}.
\end{equation}
Veamos por inducción que $\abs{a_n}\le b_n$. Tomando módulos en
\eqref{eq:recurrencia}, usando \eqref{eq:cotas-cauchy} y cambiando
$j=n-k$,
\begin{align*}
  (n+2)(n+1)\abs{a_{n+2}}
  &\le\sum_{k=0}^{n}\bigl[(n-k+1)\abs{p_k}\abs{a_{n-k+1}}+\abs{q_k}\abs{a_{n-k}}\bigr]\\
  &\le M\sum_{j=0}^{n}r^{\,j-n}\bigl[(j+1)\,b_{j+1}+b_j\bigr]
  \;\le\;(n+2)(n+1)\,b_{n+2}.
\end{align*}

\emph{Paso 3: el cociente $b_{n+2}/b_{n+1}$.} Escribiendo
\eqref{eq:mayorante} para $n$ (multiplicada por $r$) y para $n-1$, con
$n\ge1$, y restando, todos los términos de la suma salvo el de $j=n$ se
cancelan:
\[
  r\,(n+2)(n+1)\,b_{n+2}-(n+1)n\,b_{n+1}
  =M r\bigl[(n+1)b_{n+1}+b_n\bigr]+Mr^{2}b_{n+1}-Mr\,b_n ,
\]
es decir,
\begin{equation}\label{eq:cociente}
  \frac{b_{n+2}}{b_{n+1}}
  =\frac{n(n+1)+Mr\,(n+1)+Mr^{2}}{r\,(n+2)(n+1)}
  \;\xrightarrow[n\to\infty]{}\;\frac1r .
\end{equation}
(Todos los $b_n$ con $n\ge2$ son positivos: $b_2=\tfrac M2(b_0+b_1+rb_1)>0$
y \eqref{eq:cociente} propaga la positividad.)

\emph{Paso 4: convergencia.} Por el criterio del cociente,
$\sum b_n\abs{z}^{n}$ converge si $\abs{z}<r$, y por comparación
$\sum a_nz^{n}$ converge absolutamente en $\abs{z}<r$. Como $r<R$ es
arbitrario, la serie converge en $\abs{z}<R$.

\emph{Paso 5: es solución.} En $\abs{z}<R$ son lícitas la derivación término a
término (\cref{teo:series}) y el producto de Cauchy (\cref{lem:cauchy}); el
coeficiente de $z^{n}$ en $y''+py'+qy$ es exactamente el miembro izquierdo de
\eqref{eq:recurrencia}, que es nulo. Además $y(x_0)=a_0$ e $y'(x_0)=a_1$.
Esto prueba (i).

\emph{Paso 6: toda solución es una serie.} Si $\varphi$ es una solución en
$(x_0-R,x_0+R)$, la serie $y$ con $a_0=\varphi(x_0)$, $a_1=\varphi'(x_0)$ es
otra solución con los mismos datos, y por el \cref{teo:eyu}, $\varphi=y$.
Esto prueba (ii).

\emph{Paso 7: cota mediante los ceros de $P$.} Si $A$ y $B$ son enteras,
$p=A/P$ y $q=B/P$ son holomorfas en el disco
$\disco{x_0}{d}$, $d=\dist(x_0,\mathcal Z(P))$, y por el \cref{teo:radio}
sus series de Taylor convergen en él: $R_p,R_q\ge d$. Esto prueba (iii).
\end{proof}

% =========================================================================
\section{Ejemplos}\label{sec:ejemplos}
% =========================================================================

Aplicamos ahora el método a cuatro ecuaciones de complejidad creciente. En
todas ellas $x_0=0$, de modo que $z=x$. La \cref{fig:saltos} anticipa la
estructura de las recurrencias que aparecerán: cuando $p$ y $q$ son
polinomios sencillos, la recurrencia «salta» de dos en dos o de tres en
tres y desacopla los coeficientes en cadenas independientes.

% Figura: cadenas de coeficientes desacopladas (saltos de 2 y de 3)
\begin{figure}[ht]
\centering
\begin{tikzpicture}
  \tikzset{coef/.style={circle, draw=#1, fill=cajafondo, line width=0.8pt,
                        minimum size=0.82cm, inner sep=0pt, font=\small},
           coef/.default=grismedio,
           cadena/.style={-{Stealth[length=5pt,width=4.5pt]}, line width=0.7pt, #1}}
  % ---------------- (a) Salto de dos ----------------
  \begin{scope}
    \node[font=\small, text=texto, anchor=west] at (-0.5,1.45) {(a) Salto de dos: $a_{n+2}\propto a_n$};
    \foreach \i/\k in {0/0,1/2,2/4,3/6}{
      \node[coef=cian] (p\i) at (1.35*\i,0.45) {$a_{\k}$};}
    \foreach \i/\k in {0/1,1/3,2/5,3/7}{
      \node[coef=magenta] (q\i) at (1.35*\i,-0.65) {$a_{\k}$};}
    \foreach \i in {0,1,2}{
      \pgfmathtruncatemacro{\j}{\i+1}
      \draw[cadena=cian] (p\i) -- (p\j);
      \draw[cadena=magenta] (q\i) -- (q\j);}
    \node[font=\small, text=gris] at (1.35*4-0.25,0.45) {$\cdots$};
    \node[font=\small, text=gris] at (1.35*4-0.25,-0.65) {$\cdots$};
    \node[etiqgris, anchor=north west, align=left] at (-0.5,-1.3)
      {cadena par $\propto a_0$, cadena impar $\propto a_1$\\
       $y''+y=0$, \ Legendre, \ $y''+xy'+y=0$};
  \end{scope}
  % ---------------- (b) Salto de tres ----------------
  \begin{scope}[xshift=7.6cm]
    \node[font=\small, text=texto, anchor=west] at (-0.5,1.45) {(b) Salto de tres: $a_{n+2}\propto a_{n-1}$};
    \foreach \i/\k in {0/0,1/3,2/6,3/9}{
      \node[coef=cian] (r\i) at (1.35*\i,0.45) {$a_{\k}$};}
    \foreach \i/\k in {0/1,1/4,2/7,3/10}{
      \node[coef=magenta] (s\i) at (1.35*\i,-0.65) {$a_{\k}$};}
    \foreach \i/\k in {0/2,1/5,2/8,3/11}{
      \node[coef, text=gris] (t\i) at (1.35*\i,-1.75) {$a_{\k}$};}
    \foreach \i in {0,1,2}{
      \pgfmathtruncatemacro{\j}{\i+1}
      \draw[cadena=cian] (r\i) -- (r\j);
      \draw[cadena=magenta] (s\i) -- (s\j);
      \draw[cadena=grismedio, densely dashed] (t\i) -- (t\j);}
    \foreach \y in {0.45,-0.65,-1.75}{
      \node[font=\small, text=gris] at (1.35*4-0.25,\y) {$\cdots$};}
    \node[etiqgris, anchor=west] at (1.35*4+0.15,-1.75) {$=0$};
    \node[etiqgris, anchor=north west, align=left] at (-0.5,-2.4)
      {$a_2=0$ anula la tercera cadena; ecuación de Airy};
  \end{scope}
\end{tikzpicture}
\caption{Cuando $p$ y $q$ son polinomios sencillos, la recurrencia
desacopla los coeficientes en cadenas independientes. (a) Si sólo
intervienen $a_n$ y $a_{n+2}$, la cadena par depende de $a_0$ y la impar de
$a_1$. (b) En la ecuación de Airy intervienen $a_{n-1}$ y $a_{n+2}$: hay tres
cadenas, y la que arranca en $a_2=0$ es idénticamente nula.}
\label{fig:saltos}
\end{figure}


% -------------------------------------------------------------------------
\subsection{Comprobación: el oscilador armónico}\label{sec:oscilador}
% -------------------------------------------------------------------------

\begin{ejemplo}[$y''+y=0$]\label{ej:oscilador}
Aquí $p\equiv0$ y $q\equiv1$, es decir, $q_0=1$ y el resto de coeficientes
son nulos. En \eqref{eq:recurrencia} sólo sobrevive el término $k=0$ de
$q$:
\begin{equation}\label{eq:rec-oscilador}
  (n+2)(n+1)\,a_{n+2}+a_n=0
  \quad\Longrightarrow\quad
  a_{n+2}=-\frac{a_n}{(n+2)(n+1)} .
\end{equation}
La recurrencia salta de dos en dos (\cref{fig:saltos}\,(a)) y una inducción
inmediata da $a_{2m}=(-1)^{m}a_0/(2m)!$ y $a_{2m+1}=(-1)^{m}a_1/(2m+1)!$:
\begin{itemize}
  \item con $a_0=1$, $a_1=0$: $a_2=-\frac1{2!}$, $a_4=\frac1{4!},\dots$, e
        impares nulos;
  \item con $a_0=0$, $a_1=1$: $a_3=-\frac1{3!}$, $a_5=\frac1{5!},\dots$, y
        pares nulos.
\end{itemize}
Por tanto
\begin{equation}\label{eq:cos-sen}
  y_1=\sum_{m=0}^{\infty}\frac{(-1)^{m}x^{2m}}{(2m)!}=\cos x,
  \qquad
  y_2=\sum_{m=0}^{\infty}\frac{(-1)^{m}x^{2m+1}}{(2m+1)!}=\sin x .
\end{equation}
Recuperamos lo que ya sabíamos, lo cual es una buena señal. Como $P\equiv1$
no tiene ceros, el radio es $R=\infty$ (\cref{fig:coseno}).
\end{ejemplo}

% Figura: sumas parciales de la serie del coseno (R = infinito)
\begin{figure}[ht]
\centering
\begin{tikzpicture}
  \begin{axis}[oscuro, width=0.92\textwidth, height=6.0cm,
      xmin=-7, xmax=7, ymin=-2.2, ymax=2.2,
      xtick={-6.2832,-3.1416,0,3.1416,6.2832},
      xticklabels={$-2\pi$,$-\pi$,$0$,$\pi$,$2\pi$},
      ytick={-2,-1,0,1,2}, xlabel={$x$},
      samples=400, domain=-7:7, restrict y to domain=-4:4,
      legend style={at={(0.5,-0.22)}, anchor=north, legend columns=5,
                    /tikz/every even column/.append style={column sep=8pt}}]
    \addplot[c2] {1-x^2/2};
    \addplot[c3] {1-x^2/2+x^4/24};
    \addplot[c4] {1-x^2/2+x^4/24-x^6/720+x^8/40320};
    \addplot[c5] {1-x^2/2+x^4/24-x^6/720+x^8/40320-x^10/3628800+x^12/479001600};
    \addplot[blanco, line width=1.6pt] {cos(deg(x))};
    \legend{$S_2$, $S_4$, $S_8$, $S_{12}$, $\cos x$}
  \end{axis}
\end{tikzpicture}
\caption{Sumas parciales $S_{2m}(x)=\sum_{k=0}^{m}(-1)^k x^{2k}/(2k)!$ de la
solución $y_1=\cos x$ del \cref{ej:oscilador}. Como $P\equiv1$ no tiene
ceros, el radio es infinito: al aumentar el orden, la aproximación es buena
en intervalos cada vez mayores, sin barrera alguna.}
\label{fig:coseno}
\end{figure}


% -------------------------------------------------------------------------
\subsection{La ecuación de Airy}\label{sec:airy}
% -------------------------------------------------------------------------

\begin{ejemplo}[$y''-x\,y=0$]\label{ej:airy}
Ahora $p\equiv0$ y $q=-x$, es decir, $q_1=-1$ y el resto nulos. En el
sumatorio de \eqref{eq:recurrencia} sólo sobrevive $k=1$, que existe si
$n\ge1$:
\begin{equation}\label{eq:rec-airy}
  (n+2)(n+1)\,a_{n+2}-a_{n-1}=0\quad(n\ge1),\qquad a_2=0\quad(n=0).
\end{equation}
La recurrencia salta de tres en tres, $a_{n+2}=a_{n-1}/\bigl((n+2)(n+1)\bigr)$
(\cref{fig:saltos}\,(b)):
\begin{equation}\label{eq:coef-airy}
  a_3=\frac{a_0}{3\cdot2},\quad a_6=\frac{a_0}{6\cdot5\cdot3\cdot2},\qquad
  a_4=\frac{a_1}{4\cdot3},\quad a_7=\frac{a_1}{7\cdot6\cdot4\cdot3},\qquad
  a_2=a_5=a_8=\dots=0 .
\end{equation}
En general, como $(3j-1)(3j)=9j\,(j-\tfrac13)$ y $3j\,(3j+1)=9j\,(j+\tfrac13)$,
\begin{equation}\label{eq:coef-airy-gamma}
  a_{3m}=\frac{\Gamma(\tfrac23)}{9^{m}\,m!\,\Gamma(m+\tfrac23)}\,a_0,
  \qquad
  a_{3m+1}=\frac{\Gamma(\tfrac43)}{9^{m}\,m!\,\Gamma(m+\tfrac43)}\,a_1 ,
\end{equation}
y el sistema fundamental es
\begin{equation}\label{eq:airy-y1y2}
  y_1=1+\frac{x^{3}}{6}+\frac{x^{6}}{180}+\cdots,\qquad
  y_2=x+\frac{x^{4}}{12}+\frac{x^{7}}{504}+\cdots ,
\end{equation}
con radio $R=\infty$. Estas funciones \emph{no son elementales}: el método de
series \emph{define} funciones nuevas. Las funciones de Airy de uso habitual
son las combinaciones \cite[§9.4]{DLMF}
\begin{equation}\label{eq:Ai-Bi}
  \Ai(x)=c_1\,y_1(x)-c_2\,y_2(x),\qquad
  \Bi(x)=\sqrt3\,\bigl(c_1\,y_1(x)+c_2\,y_2(x)\bigr),
\end{equation}
con $c_1=\Ai(0)=3^{-2/3}/\Gamma(\tfrac23)\approx0{,}355028$ y
$c_2=-\Ai'(0)=3^{-1/3}/\Gamma(\tfrac13)\approx0{,}258819$.
\end{ejemplo}

\begin{observacion}[Interpretación cualitativa]\label{obs:airy}
Para $x<0$ la ecuación es $y''=-\abs{x}\,y$: un oscilador de «frecuencia»
$\sqrt{\abs{x}}$ creciente, y las soluciones oscilan cada vez más deprisa
con amplitud decreciente. Para $x>0$ es $y''=x\,y$ y las soluciones crecen o
decrecen exponencialmente; $\Ai$ es la única (salvo múltiplos) que decae.
La ecuación de Airy es la ecuación de Schrödinger en un potencial lineal y
describe el comportamiento de las funciones de onda cerca de un punto de
retroceso clásico \cite[§10.2]{BENDER}.
\end{observacion}

% Figura: soluciones de la ecuación de Airy (datos: scripts/generar_datos.py)
\begin{figure}[ht]
\centering
\begin{tikzpicture}
  \begin{groupplot}[group style={group size=1 by 2, vertical sep=1.5cm},
      oscuro, width=0.92\textwidth, height=5.3cm,
      xmin=-9, xmax=2.4, xtick={-8,-6,-4,-2,0,2}, unbounded coords=jump]
    % ---------- (a) sistema fundamental ----------
    \nextgroupplot[ymin=-1.6, ymax=3.2, ytick={-1,0,1,2,3},
        title={(a) Sistema fundamental $y_1$, $y_2$ (series \eqref{eq:airy-y1y2})},
        legend pos=north west]
      \fill[grisoscuro, opacity=0.45] (axis cs:0,-1.6) rectangle (axis cs:2.4,3.2);
      \addplot[c1, restrict y to domain=-10:10] table[x=x, y=y1] {datos/airy.dat};
      \addplot[c2, restrict y to domain=-10:10] table[x=x, y=y2] {datos/airy.dat};
      \legend{$y_1$, $y_2$}
      \node[font=\scriptsize, text=gris, anchor=south] at (axis cs:-6.6,-1.58) {región oscilatoria ($x<0$)};
      \node[font=\scriptsize, text=gris, anchor=south] at (axis cs:1.2,-1.55) {crecimiento};
    % ---------- (b) funciones de Airy ----------
    \nextgroupplot[ymin=-0.65, ymax=1.6, ytick={-0.5,0,0.5,1,1.5},
        xlabel={$x$},
        title={(b) Funciones de Airy $\Ai$ y $\Bi$, combinaciones \eqref{eq:Ai-Bi}},
        legend pos=north west]
      \fill[grisoscuro, opacity=0.45] (axis cs:0,-0.65) rectangle (axis cs:2.4,1.6);
      \addplot[c3, restrict y to domain=-10:10] table[x=x,
          y expr={0.355028054*\thisrow{y1}-0.258819404*\thisrow{y2}}] {datos/airy.dat};
      \addplot[c4, restrict y to domain=-10:10] table[x=x,
          y expr={1.732050808*(0.355028054*\thisrow{y1}+0.258819404*\thisrow{y2})}] {datos/airy.dat};
      \legend{$\Ai$, $\Bi$}
  \end{groupplot}
\end{tikzpicture}
\caption{Ecuación de Airy $y''=x\,y$. Para $x<0$ las soluciones oscilan con
frecuencia creciente y amplitud decreciente; para $x>0$ (zona sombreada)
crecen exponencialmente, salvo $\Ai$, que decae. Las curvas se han obtenido
sumando directamente las series \eqref{eq:airy-y1y2}, que convergen en todo
$\R$.}
\label{fig:airy}
\end{figure}


% -------------------------------------------------------------------------
\subsection{Un problema de valores iniciales}\label{sec:pvi}
% -------------------------------------------------------------------------

\begin{ejemplo}[$y''+x\,y'+y=0$, $y(0)=y_0$, $y'(0)=v_0$]\label{ej:pvi}
Este problema procede de \cite{TEMAV}. Ahora $p=x$ y $q\equiv1$, es decir,
$p_1=1$ y $q_0=1$. En \eqref{eq:recurrencia} sobreviven el término $k=1$ de
$p$, que vale $n\,p_1\,a_n$, y el término $k=0$ de $q$:
\begin{equation}\label{eq:rec-pvi}
  (n+2)(n+1)\,a_{n+2}+(n+1)\,a_n=0
  \quad\Longrightarrow\quad
  a_{n+2}=-\frac{a_n}{n+2},\qquad n\ge0 .
\end{equation}
De nuevo la recurrencia salta de dos en dos. Con la notación del doble
factorial, $(2i)!!=2\cdot4\cdots(2i)=2^{i}\,i!$ y
$(2i+1)!!=1\cdot3\cdots(2i+1)$,
\begin{equation}\label{eq:coef-pvi}
  a_{2i}=\frac{(-1)^{i}}{(2i)!!}\,a_0,\qquad
  a_{2i+1}=\frac{(-1)^{i}}{(2i+1)!!}\,a_1 ,
\end{equation}
de donde
\begin{equation}\label{eq:pvi-y1y2}
  y_1=\sum_{i=0}^{\infty}\frac{(-1)^{i}}{(2i)!!}\,x^{2i}
     =\sum_{i=0}^{\infty}\frac{1}{i!}\Bigl(-\frac{x^{2}}{2}\Bigr)^{i}
     =\eul^{-x^{2}/2},
  \qquad
  y_2=\sum_{i=0}^{\infty}\frac{(-1)^{i}}{(2i+1)!!}\,x^{2i+1}.
\end{equation}
Imponiendo las condiciones iniciales, $a_0=y_0$ y $a_1=v_0$, y la solución
del problema es $y=y_0\,y_1+v_0\,y_2$, con radio $R=\infty$.
\end{ejemplo}

\begin{observacion}[Comprobación por integración directa]\label{obs:pvi-exacta}
La ecuación es exacta: $y''+x\,y'+y=\bigl(y'+x\,y\bigr)'$. Integrando una vez
con los datos iniciales, $y'+x\,y=v_0$, ecuación lineal de primer orden con
factor integrante $\eul^{x^{2}/2}$:
\[
  y(x)=\eul^{-x^{2}/2}\Bigl(y_0+v_0\int_0^{x}\eul^{t^{2}/2}\,\dif t\Bigr).
\]
Por la unicidad del \cref{teo:eyu}, $y_1=\eul^{-x^{2}/2}$ (confirmando
\eqref{eq:pvi-y1y2}) e
$y_2=\eul^{-x^{2}/2}\int_0^{x}\eul^{t^{2}/2}\,\dif t$, que es, salvo
reescalado, la función de Dawson: $y_2(x)=\sqrt2\,F(x/\sqrt2)$ con
$F(u)=\eul^{-u^{2}}\int_0^{u}\eul^{s^{2}}\dif s$ \cite[§7.2]{DLMF}. El método
de series ha reconstruido, sin integrar, una gaussiana y una función no
elemental.
\end{observacion}

% Figura: soluciones de y'' + x y' + y = 0 (datos: scripts/generar_datos.py)
\begin{figure}[ht]
\centering
\begin{tikzpicture}
  \begin{axis}[oscuro, width=0.92\textwidth, height=5.6cm,
      xmin=-5, xmax=5, ymin=-1.15, ymax=1.65,
      xtick={-4,-2,0,2,4}, ytick={-1,-0.5,0,0.5,1,1.5}, xlabel={$x$},
      legend style={at={(0.985,0.04)}, anchor=south east}]
    \addplot[c1] table[x=x, y=y1] {datos/ejemplo_gauss.dat};
    \addplot[c2] table[x=x, y=y2] {datos/ejemplo_gauss.dat};
    \addplot[c3, densely dashed] table[x=x, y expr={\thisrow{y1}+\thisrow{y2}}] {datos/ejemplo_gauss.dat};
    \legend{$y_1=\eul^{-x^2/2}$, $y_2=\eul^{-x^2/2}\!\int_0^x\eul^{t^2/2}\dif t$,
            $y_1+y_2$ \ ($y_0=v_0=1$)}
  \end{axis}
\end{tikzpicture}
\caption{Sistema fundamental del \cref{ej:pvi}: la gaussiana $y_1$ (par) y
la función de tipo Dawson $y_2$ (impar, con $y_2\sim1/x$ para $\abs{x}\to\infty$),
junto con la solución del problema de valores iniciales con $y_0=v_0=1$.}
\label{fig:gauss}
\end{figure}


% -------------------------------------------------------------------------
\subsection{La ecuación de Legendre}\label{sec:legendre}
% -------------------------------------------------------------------------
\paragraph{Truco práctico: no dividir por $P$}
Si $P$ es un polinomio, suele ser más cómodo sustituir la serie
directamente en \eqref{eq:edo}, sin pasar a la forma normal: así se evita
desarrollar los cocientes $A/P$ y $B/P$. Las multiplicaciones por potencias
de $x$ sólo desplazan índices: por ejemplo,
$x^{2}\,y''=\sum_{n\ge0}n(n-1)\,a_n\,x^{n}$ y $x\,y'=\sum_{n\ge0}n\,a_n\,x^{n}$
(los términos añadidos con $n=0,1$ son nulos).

\begin{ejemplo}[Legendre en $x_0=0$]\label{ej:legendre}
Consideremos
\begin{equation}\label{eq:legendre}
  (1-x^{2})\,y''-2x\,y'+\ell(\ell+1)\,y=0 .
\end{equation}
Sustituyendo y agrupando potencias $x^{n}$ ---el término $-x^{2}y''$ aporta
$-n(n-1)\,a_n$ y $-2x\,y'$ aporta $-2n\,a_n$---,
\begin{equation}\label{eq:rec-legendre-0}
  (n+2)(n+1)\,a_{n+2}-\bigl[n(n+1)-\ell(\ell+1)\bigr]\,a_n=0 ,
\end{equation}
y, factorizando $n(n+1)-\ell(\ell+1)=(n-\ell)(n+\ell+1)$,
\begin{equation}\label{eq:rec-legendre}
  \boxed{\;a_{n+2}=\frac{(n-\ell)(n+\ell+1)}{(n+2)(n+1)}\,a_n ,\qquad n\ge0.\;}
\end{equation}
El radio de convergencia garantizado es $R=1$, distancia de $0$ a los ceros
$x=\pm1$ de $P$. Para una serie que no se corte es exactamente $1$: como
$a_{n+2}/a_n\to1$, el criterio del cociente aplicado a la serie en $x^{2}$
da divergencia para $\abs{x}>1$.
\end{ejemplo}

\begin{proposicion}[Polinomios de Legendre]\label{prop:polinomios-legendre}
Si $\ell\in\N$, la serie de la misma paridad que $\ell$ se trunca: el
numerador de \eqref{eq:rec-legendre} se anula para $n=\ell$, luego
$a_{\ell+2}=a_{\ell+4}=\dots=0$, y la solución correspondiente es un
polinomio de grado $\ell$. Normalizado por $P_\ell(1)=1$, es el
\emph{polinomio de Legendre} $P_\ell$:
\[
  P_0=1,\quad P_1=x,\quad P_2=\tfrac12(3x^{2}-1),\quad
  P_3=\tfrac12(5x^{3}-3x),\quad P_4=\tfrac18(35x^{4}-30x^{2}+3).
\]
La serie de la otra paridad no se corta y tiene radio exactamente $1$.
\end{proposicion}

\begin{proof}
El truncamiento es inmediato a partir de \eqref{eq:rec-legendre}. Por
ejemplo, para $\ell=4$: $a_2=\frac{(-4)\cdot5}{2}\,a_0=-10\,a_0$,
$a_4=\frac{(-2)\cdot7}{12}\,a_2=\frac{35}{3}\,a_0$ y $a_6=0$; imponiendo
$a_0\bigl(1-10+\frac{35}3\bigr)=1$ se obtiene $a_0=\frac38$ y la expresión
de $P_4$. Si la otra serie se cortara, ambas soluciones serían polinomios y
estarían acotadas en $[-1,1]$; el apartado siguiente muestra (para $\ell=1$)
que no es así, y el caso general se trata en \cite[§15.6]{ARFKEN}.
\end{proof}

% Figura: polinomios de Legendre y serie no truncada para l = 1
\begin{figure}[ht]
\centering
\begin{tikzpicture}
  \begin{groupplot}[group style={group size=2 by 1, horizontal sep=1.3cm},
      oscuro, width=0.5\textwidth, height=6.2cm, xmin=-1, xmax=1,
      xtick={-1,-0.5,0,0.5,1}, xlabel={$x$}]
    % ---------- (a) P_0 ... P_4 ----------
    \nextgroupplot[ymin=-1.1, ymax=1.1, ytick={-1,-0.5,0,0.5,1},
        title={(a) Polinomios $P_\ell$ (series truncadas)},
        samples=200, domain=-1:1,
        legend style={at={(0.5,-0.24)}, anchor=north, legend columns=5,
                      /tikz/every even column/.append style={column sep=3pt}}]
      \addplot[line width=1.15pt, gris] {1};
      \addplot[c1] {x};
      \addplot[c2] {0.5*(3*x^2-1)};
      \addplot[c3] {0.5*(5*x^3-3*x)};
      \addplot[c4] {(35*x^4-30*x^2+3)/8};
      \legend{$P_0$, $P_1$, $P_2$, $P_3$, $P_4$}
    % ---------- (b) y1 = 1 - x artanh x ----------
    \nextgroupplot[ymin=-2.6, ymax=1.25, ytick={-2,-1,0,1},
        title={(b) Serie par para $\ell=1$: no se corta},
        legend style={at={(0.5,-0.24)}, anchor=north, legend columns=5,
                      /tikz/every even column/.append style={column sep=3pt}}]
      \draw[guia] (axis cs:-1,-2.6) -- (axis cs:-1,1.25);
      \draw[guia] (axis cs:1,-2.6) -- (axis cs:1,1.25);
      \addplot[c2] table[x=x, y=S1] {datos/legendre_l1_parciales.dat};
      \addplot[c3] table[x=x, y=S2] {datos/legendre_l1_parciales.dat};
      \addplot[c4] table[x=x, y=S4] {datos/legendre_l1_parciales.dat};
      \addplot[c5] table[x=x, y=S8] {datos/legendre_l1_parciales.dat};
      \addplot[blanco, line width=1.6pt, samples=400, domain=-0.997:0.997]
        {1-0.5*x*ln((1+x)/(1-x))};
      \legend{$S_1$, $S_2$, $S_4$, $S_8$, $y_1$}
  \end{groupplot}
\end{tikzpicture}
\caption{(a) Para $\ell\in\N$ la serie de la paridad de $\ell$ se corta y da
el polinomio de Legendre $P_\ell$, regular en todo $[-1,1]$ (y, de hecho, en
todo $\R$). (b) Para $\ell=1$, la serie par
$y_1=1-x\artanh x$ de \eqref{eq:legendre-q1} y sus sumas parciales
$S_M=1-\sum_{k=1}^{M}x^{2k}/(2k-1)$: la convergencia se degrada al
acercarse a $x=\pm1$, donde $y_1$ diverge logarítmicamente. El radio $R=1$
del \cref{teo:fuchs} se alcanza.}
\label{fig:legendre}
\end{figure}


\begin{ejemplo}[Legendre con $\ell=1$]\label{ej:legendre-1}
Para $(1-x^{2})\,y''-2x\,y'+2y=0$ \cite{TEMAV}, la recurrencia
\eqref{eq:rec-legendre} se reduce a
\begin{equation}\label{eq:rec-legendre-1}
  a_{n+2}=\frac{(n-1)(n+2)}{(n+2)(n+1)}\,a_n=\frac{n-1}{n+1}\,a_n .
\end{equation}
\begin{itemize}
  \item \emph{Serie impar.} $a_3=0\cdot a_1=0$, y con él todos los impares
        posteriores: la serie impar se reduce a $y_2=x=P_1(x)$.
  \item \emph{Serie par.} $a_2=-a_0$, $a_4=\tfrac13a_2=-\tfrac13a_0$,
        $a_6=\tfrac35a_4=-\tfrac15a_0$ y, por inducción,
        $a_{2k}=-a_0/(2k-1)$ para $k\ge1$. Por tanto
        \begin{equation}\label{eq:legendre-q1}
          y_1=1-\sum_{k=1}^{\infty}\frac{x^{2k}}{2k-1}
             =1-x\,\artanh x
             =1-\frac{x}{2}\,\ln\frac{1+x}{1-x},\qquad\abs{x}<1 .
        \end{equation}
\end{itemize}
La función $y_1$ (que es $-Q_1$, con $Q_1$ la función de Legendre de segunda
especie) diverge logarítmicamente en $x=\pm1$: la cota del \cref{teo:fuchs}
es \emph{óptima} en este caso. Obsérvese también que la ecuación
$(1-x^{2})\,y''-2x\,y'+\tfrac34y=0$ del \cref{ej:ordinario} es la de Legendre
con $\ell(\ell+1)=\tfrac34$, es decir, $\ell=\tfrac12\notin\N$: ninguna de sus
dos series se corta y ambas tienen radio exactamente $1$.
\end{ejemplo}

\begin{observacion}[Origen físico y «cuantización»]\label{obs:legendre-fisica}
La ecuación de Legendre aparece al separar variables en la ecuación de
Laplace en coordenadas esféricas, con $x=\cos\theta$. Los puntos $x=\pm1$
corresponden a los polos $\theta=0,\pi$, donde la solución física debe ser
finita. Por \eqref{eq:legendre-q1} y la \cref{prop:polinomios-legendre}, eso
sólo es posible si una serie se trunca, lo que exige $\ell\in\N$: la
condición de regularidad \emph{cuantiza} el parámetro, exactamente como la
condición de normalizabilidad cuantiza la energía en la ecuación de Hermite
del oscilador cuántico \cite{TEMAV}, \cite[§13.1]{ARFKEN}.
\end{observacion}

% =========================================================================
\section{Resumen del método y perspectivas}\label{sec:resumen}
% =========================================================================

\begin{metodo}[Receta para un punto ordinario $x_0$]
\begin{enumerate}[label=\textcolor{ambar}{\bfseries\arabic*.}, leftmargin=1.6em]
  \item Comprobar que $P(x_0)\neq0$ y estimar
        $R\ge\dist\bigl(x_0,\mathcal Z(P)\bigr)$, la distancia al cero
        \emph{complejo} de $P$ más cercano (\cref{teo:fuchs}).
  \item Proponer $y=\sum_{n\ge0}a_nz^{n}$, con $z=x-x_0$, y derivar término
        a término.
  \item Sustituir ---en la forma normal o, si $P$ es un polinomio,
        directamente en \eqref{eq:edo}---, desplazar índices para que todo
        vaya en $z^{n}$ y usar el producto de Cauchy.
  \item Igualar a cero el coeficiente de cada $z^{n}$
        (\cref{cor:unicidad}): se obtiene la recurrencia para $a_{n+2}$.
  \item Dejar libres $a_0=y(x_0)$ y $a_1=y'(x_0)$: $y_1$ corresponde a
        $(a_0,a_1)=(1,0)$ e $y_2$ a $(0,1)$; la solución general es
        $y=a_0y_1+a_1y_2$.
\end{enumerate}
\end{metodo}

La \cref{fig:flujo} recoge el mismo algoritmo en forma de diagrama, y la
\cref{tab:resumen} compara los cuatro ejemplos de la \cref{sec:ejemplos}.

\begin{table}[!htb]
\centering
\caption{Resumen de los ejemplos resueltos en $x_0=0$.}
\label{tab:resumen}
\renewcommand{\arraystretch}{2.25}
\small
\begin{tabular}{@{}l l l c l@{}}
\toprule
\textcolor{cian}{\bfseries Ecuación} & \textcolor{cian}{\bfseries Recurrencia} &
\textcolor{cian}{\bfseries Salto} & \textcolor{cian}{\bfseries $R$} &
\textcolor{cian}{\bfseries Soluciones}\\
\midrule
$y''+y=0$ & $a_{n+2}=-\dfrac{a_n}{(n+2)(n+1)}$ & $2$ & $\infty$ & $\cos x$, $\sin x$\\
$y''-xy=0$ & $a_{n+2}=\dfrac{a_{n-1}}{(n+2)(n+1)}$ & $3$ & $\infty$ & $\Ai$, $\Bi$ (no elementales)\\
$y''+xy'+y=0$ & $a_{n+2}=-\dfrac{a_n}{n+2}$ & $2$ & $\infty$ & $\eul^{-x^2/2}$, tipo Dawson\\
Legendre & $a_{n+2}=\dfrac{(n-\ell)(n+\ell+1)}{(n+2)(n+1)}\,a_n$ & $2$ & $1$ & $P_\ell$ si $\ell\in\N$; otra con $R=1$\\
\bottomrule
\end{tabular}
\end{table}

% Figura: diagrama de flujo del método
\begin{figure}[!htb]
\centering
\begin{tikzpicture}[node distance=0.55cm]
  \tikzset{paso/.style={nodo, text width=7.4cm, align=left, minimum height=0.85cm},
           num/.style={font=\sffamily\bfseries\small, text=ambar}}
  \node[nodocian, align=center, text width=5.4cm] (ini)
    {$P(x)\,y''+A(x)\,y'+B(x)\,y=0$\\[1pt]\footnotesize\color{gris} y un punto $x_0$};
  \node[diamond, aspect=2.4, draw=violeta, fill=cajafondo, line width=0.6pt,
        inner sep=1.5pt, font=\small, below=of ini] (dec) {¿$P(x_0)\neq0$?};
  \node[paso, below=0.75cm of dec] (p1)
    {\textcolor{ambar}{\bfseries 1.}\ \ $R\ge\dist\bigl(x_0,\mathcal Z(P)\bigr)$\quad(ceros complejos)};
  \node[paso, below=of p1] (p2)
    {\textcolor{ambar}{\bfseries 2.}\ \ $y=\sum a_n z^{n}$; derivar término a término};
  \node[paso, below=of p2] (p3)
    {\textcolor{ambar}{\bfseries 3.}\ \ desplazar índices y producto de Cauchy};
  \node[paso, below=of p3] (p4)
    {\textcolor{ambar}{\bfseries 4.}\ \ coeficiente de $z^{n}$ nulo $\Rightarrow$ recurrencia};
  \node[paso, below=of p4] (p5)
    {\textcolor{ambar}{\bfseries 5.}\ \ $(1,0)\mapsto y_1$, \ $(0,1)\mapsto y_2$; \ $y=a_0y_1+a_1y_2$};
  \node[nodo, draw=rojo, align=center, text width=3.6cm, right=1.6cm of dec] (sing)
    {\textcolor{rojo}{\bfseries Punto singular}\\[1pt]
     \footnotesize\color{gris} serie de Frobenius\\(Parte II)};

  \draw[flecha] (ini) -- (dec);
  \draw[flecha] (dec) -- node[right=2pt, etiqgris] {sí} (p1);
  \draw[flecha] (dec) -- node[above=2pt, etiqgris] {no} (sing);
  \foreach \a/\b in {p1/p2, p2/p3, p3/p4, p4/p5}{\draw[flecha] (\a) -- (\b);}
\end{tikzpicture}
\caption{Diagrama de flujo del método de la serie entera.}
\label{fig:flujo}
\end{figure}



\paragraph{Perspectivas}
El método de la serie entera falla cuando $P(x_0)=0$: en un punto singular
la solución puede no ser analítica (piénsese en $y=x^{r}$ con $r\notin\N$,
o en $y=\ln x$). La segunda parte del tema \cite[cap.~V]{TEMAV} estudia los
puntos \emph{singulares regulares} (\cref{obs:avance}) mediante series de
Frobenius $y=(x-x_0)^{r}\sum_n a_n(x-x_0)^n$, cuyo exponente $r$ fija la
\emph{ecuación indicial}; el caso de raíces que difieren en un entero
conduce a las \emph{series logarítmicas}. El cambio $\bar x=1/x$ permite,
además, analizar el comportamiento de las soluciones en un entorno del
infinito. Finalmente, las ecuaciones clásicas de la física ---Hermite,
Chebyshev, Legendre--- comparten la estructura de la \cref{sec:legendre}:
una recurrencia de salto dos cuya serie sólo se trunca para valores
discretos del parámetro, lo que da lugar a familias de polinomios
ortogonales y a la cuantización de magnitudes físicas
\cite{ARFKEN,RILEY}.


% =========================================================================
%                              BIBLIOGRAFÍA
% =========================================================================
\FloatBarrier
\printbibliography[heading=bibnumbered, title={Referencias}]

\end{document}
